% આ ભાગનો પૂર્ણ સંપાદનીય પાઠ છે. શૈલી, ફોન્ટ, ચિત્રો અને પ્રકરણસંદર્ભ માટે સંપૂર્ણ સ્રોત ZIP તથા reader/full-reader.tex જુઓ. આ એકલી સ્વતંત્ર કમ્પાઇલ થતી મુખ્ય ફાઇલ નથી.

% BEGIN OLP-0491
\olpart{int}{સ્ફુરણાવાદી તર્કશાસ્ત્ર}
\phantomsection\label{guunit:OLP-0491}


\begingroup
\makeatletter\def\ol@sectioncs{section}\makeatother

% BEGIN OLP-0492
\olchapter{int}{int}{પરિચય}
\olchapterid{int}{int}
\phantomsection\label{guunit:OLP-0492}


\begingroup
\makeatletter\def\ol@sectioncs{section}\makeatother

% BEGIN OLP-0493
\olfileid{int}{int}{cr}

\section{રચનાત્મક તર્કવિચાર}
\phantomsection\label{guunit:OLP-0493}


મોડલ સંક્રિયકો કે દ્વિતીય-ક્રમ પરિમાણકો ઉમેરીને પ્રશિષ્ટ તર્કનો
વિસ્તાર કરવાથી ઊલટું, સ્ફુરણાવાદી તર્ક પ્રશિષ્ટ તર્કને મર્યાદિત
કરે છે; એ અર્થમાં તે ``અપ્રશિષ્ટ'' છે. પ્રશિષ્ટ તર્ક અનેક રીતે
\emph{અરચનાત્મક} છે. સ્ફુરણાવાદી તર્કનો આશય રચનાત્મક ગણિતમાં
જોવા મળતા વધુ ``રચનાત્મક'' પ્રકારના તર્કવિચારનું નિદર્શન કરવાનો
છે. નીચેના દાખલા તેની કેટલીક પાયાની પ્રેરણાઓ સ્પષ્ટ કરે છે.

ધારો કે કોઈ દાવો કરે કે તેણે એવી પ્રાકૃતિક સંખ્યા~$n$ નક્કી
કરી છે કે જો $n$ સમ હોય તો રીમાન પરિકલ્પના સત્ય છે, અને જો
$n$ વિષમ હોય તો રીમાન પરિકલ્પના અસત્ય છે. બહુ સારા સમાચાર!{}
રીમાન પરિકલ્પના સત્ય છે કે નહીં એ ગણિતના મોટા વણઉકેલ્યા
પ્રશ્નોમાંનો એક છે; અહીં તો જાણે સમસ્યા ગણતરીના પ્રશ્નમાં,
એટલે કોઈ ચોક્કસ સંખ્યા સમ છે કે નહીં તે નક્કી કરવામાં, ઘટી
ગઈ છે.

તો~$n$નું એ જાદુઈ મૂલ્ય શું છે? તેઓ તેનું વર્ણન આ રીતે કરે છે:
રીમાન પરિકલ્પના સત્ય હોય તો $n$ બરાબર $2$ અને નહીં તો $3$
થાય એવી પ્રાકૃતિક સંખ્યા.

ગુસ્સામાં તમે તમારા પૈસા પાછા માગો છો. પ્રશિષ્ટ દૃષ્ટિએ ઉપરનું
વર્ણન ખરેખર $n$નું એકમાત્ર મૂલ્ય નક્કી કરે છે; પરંતુ તમને
ખરેખર તો \emph{સ્પષ્ટ રીતે} આપેલું $n$નું મૂલ્ય જોઈએ છે.

બીજો, કદાચ થોડો ઓછો કૃત્રિમ, દાખલો જોઈએ. આપણને ખબર છે કે
અસંમેય સંખ્યાને સંમેય ઘાતે ચઢાવવાથી સંમેય પરિણામ મળી શકે છે.
દાખલા તરીકે, $\sqrt{2}^2 = 2$. અસંમેય સંખ્યાને
\emph{અસંમેય} ઘાતે ચઢાવીને સંમેય પરિણામ મેળવી શકાય છે કે
નહીં તે ઓછું સ્પષ્ટ છે. નીચેનો પ્રમેય તેનો હકારાત્મક જવાબ
આપે છે:

\begin{thm}
એવી અસંમેય સંખ્યાઓ $a$ અને $b$ છે કે $a^b$ સંમેય છે.
\end{thm}

\begin{proof}
$\sqrt{2}^{\sqrt{2}}$ને વિચારો. તે સંમેય હોય, તો આપણું કામ
થઈ ગયું: $a = b = \sqrt{2}$ લઈ શકીએ. નહીં તો તે અસંમેય છે.
પછી
\[
(\sqrt{2}^{\sqrt{2}})^{\sqrt{2}} = \sqrt{2}^{\sqrt{2} \cdot
  \sqrt{2}} = \sqrt{2}^2 = 2,
\]
જે સંમેય છે. તેથી આ કિસ્સામાં $a$ને
$\sqrt{2}^{\sqrt{2}}$ અને $b$ને~$\sqrt 2$ લો.
\end{proof}

શું આ પ્રામાણ્ય સાબિતી છે? મોટા ભાગના ગણિતજ્ઞોને એમ લાગે છે.
પણ અહીં ફરી કંઈક થોડું અસંતોષકારક છે: ચોક્કસ ગુણધર્મ ધરાવતા
વાસ્તવિક સંખ્યાઓના જોડાનું અસ્તિત્વ આપણે સાબિત કર્યું છે,
પરંતુ તે સંખ્યાઓનો \emph{કયો} જોડો છે તે કહી શકતા નથી. એ જ
પરિણામ એવી રીતે પણ સાબિત કરી શકાય છે કે જોડો $a$, $b$
સાબિતીમાં જ \emph{આપેલો} હોય: $a = \sqrt{3}$ અને
$b = \log_3 4$ લો. પછી
\[
a^b = \sqrt{3}^{\log_3 4} = 3^{1/2 \cdot \log_3 4} = (3^{\log_3
  4})^{1/2} = 4^{1/2}= 2,
\]
કારણ કે $3^{\log_3 x} = x$.

સ્ફુરણાવાદી તર્કમાં પહેલા પુરાવા જેવી ચાલ માન્ય ન હોય એવા
તર્કવિચારનું નિદર્શન કરવાનો આશય છે. $!A(x)$ સંતોષતા કોઈ~$x$નું
અસ્તિત્વ સાબિત કરવાનો અર્થ, બીજા પુરાવાની જેમ, ચોક્કસ~$x$
અને તે~$!A$ સંતોષે છે તેની સાબિતી આપવાનો છે. $!A$ અથવા
$!B$ સત્ય છે તે સાબિત કરવા માટે બંનેમાંથી એકને સાબિત કરી
શકવું જરૂરી છે.

વિધિવત્ રીતે કહીએ તો, પ્રશિષ્ટ તર્કના નિષ્પત્તિ તંત્રને
ચોક્કસ રીતે મર્યાદિત કરવાથી સ્ફુરણાવાદી તર્ક મળે છે. ગાણિતિક
દૃષ્ટિએ આ માત્ર વિધિવત્ નિષ્પત્તિ-તંત્રો છે; પરંતુ અગાઉ
નોંધ્યું તેમ, તેમનો આશય એક પ્રકારના ગાણિતિક તર્કવિચારનું
નિદર્શન કરવાનો છે. આને ગણિત વિશેના ચોક્કસ તત્ત્વચિંતનાત્મક
દૃષ્ટિકોણથી ન્યાયસંગત તર્કવિચાર (જેમ કે બ્રાઉરનો સ્ફુરણાવાદ),
અથવા વધુ ``મૂર્ત'' અને સંતોષકારક ગાણિતિક તર્કવિચાર (બિશપના
રચનાવાદની દિશામાં) માની શકાય. આ વિધિવત્ વર્ણન અનૌપચારિક
પ્રેરણાને પકડે છે કે નહીં તે અંગે પણ દલીલ થઈ શકે. પરંતુ
આપણું તત્ત્વચિંતનાત્મક વલણ જે હોય તે, સ્ફુરણાવાદી તર્કનો
વિધિવત્ રીતે રજૂ થયેલા તર્ક તરીકે અભ્યાસ કરી શકીએ છીએ;
અને કારણ ગમે તે હોય, ઘણા ગાણિતિક તર્કશાસ્ત્રીઓને તેનો
અભ્યાસ રસપ્રદ લાગે છે.
% END OLP-0493
\endgroup


\begingroup
\makeatletter\def\ol@sectioncs{section}\makeatother

% BEGIN OLP-0494
\olfileid{int}{int}{syn}

\olsection{સ્ફુરણાવાદી તર્કની વાક્યરચના}
\phantomsection\label{guunit:OLP-0494}


સ્ફુરણાવાદી તર્કની વાક્યરચના વિધાનાત્મક તર્ક જેવી જ છે.
પ્રશિષ્ટ વિધાનાત્મક તર્કમાં કેટલાક સંયોજકોને બીજાં સંયોજકો
વડે વ્યાખ્યાયિત કરી શકાય છે. દાખલા તરીકે, $!A \lif !B$ને
$\lnot !A \lor !B$ વડે, અથવા $!A \lor !B$ને
$\lnot(\lnot !A \land \lnot !B)$ વડે વ્યાખ્યાયિત કરી શકીએ.
તેથી પ્રશિષ્ટ તર્કની રજૂઆતોમાં કેટલાક સંયોજકો આવી વ્યાખ્યાઓના
સંક્ષેપ તરીકે રજૂ થાય છે. સ્ફુરણાવાદી તર્કમાં એવું નથી, બે
અપવાદો સિવાય: $\lnot !A$ને $!A \lif \lfalse$ના સંક્ષેપ
તરીકે વ્યાખ્યાયિત કરી શકાય છે---અને ઘણી વાર કરાય છે. તો,
અલબત્ત, $\lfalse$ પોતે વ્યાખ્યાયિત સંક્ષેપ ન હોવું જોઈએ!{}
વળી, પ્રશિષ્ટ તર્કની જેમ $!A \liff !B$ને
$(!A \lif !B) \land (!B \lif !A)$ તરીકે વ્યાખ્યાયિત કરી
શકાય છે.

વિધાનાત્મક સ્ફુરણાવાદી તર્કનાં સૂત્રો
\emph{પ્રતિજ્ઞાપન ચલો} અને વિધાનાત્મક
અચલાંક~$\lfalse$માંથી \emph{તાર્કિક સંયોજકો} વડે રચાય છે.
આપણી પાસે છે:

\begin{enumerate}
\item પ્રતિજ્ઞાપન ચલોનો ગણનીય ગણ~$\PVar$:
  $\Obj p_0$,
  $\Obj p_1$, \dots
  \item અસત્ય માટેનો વિધાનાત્મક અચલાંક~$\lfalse$.
\item તાર્કિક સંયોજકો: $\land$ (સંયોજન), $\lor$ (વિયોજન),
  $\lif$ (શરતી સંયોજક).
\item વિરામચિહ્નો: (, ), અને અલ્પવિરામ.
\end{enumerate}

\begin{defn}[સૂત્ર]
\ollabel{defn:formulas} વિધાનાત્મક સ્ફુરણાવાદી તર્કનાં
\emph{સૂત્રો}નો ગણ~$\Frm[L_0]$ નીચે મુજબ આગમનાત્મક
રીતે વ્યાખ્યાયિત થાય છે:
\begin{enumerate}
\item $\lfalse$ પરમાણુક સૂત્ર છે.

\item દરેક પ્રતિજ્ઞાપન ચલ~$\Obj p_i$ પરમાણુક
  સૂત્ર છે.

\item જો $!A$ અને $!B$ એ સૂત્રો હોય, તો $(!A \land
  !B)$ પણ સૂત્ર છે.

\item જો $!A$ અને $!B$ એ સૂત્રો હોય, તો $(!A \lor !B)$
  પણ સૂત્ર છે.

\item જો $!A$ અને $!B$ એ સૂત્રો હોય, તો $(!A \lif !B)$
  પણ સૂત્ર છે.

\item આ સિવાય બીજું કશું સૂત્ર નથી.
\end{enumerate}
\end{defn}

ઉપર જણાવેલા મૂળભૂત સંયોજકો ઉપરાંત, નીચેના \emph{વ્યાખ્યાયિત}
સંકેતો પણ વાપરીએ છીએ: $\lnot$ (નિષેધ) અને $\liff$
(દ્વિમુખી પ્રેરણ). આ વ્યાખ્યાયિત સંક્રિયકો વડે બનેલાં
સૂત્રોને નીચે મુજબ સમજવાનાં છે:
\begin{enumerate}
\item $\lnot !A$ એ $!A \lif \lfalse$નો સંક્ષેપ છે.

\item $!A \liff !B$ એ $(!A \lif !B) \land (!B
  \lif !A)$નો સંક્ષેપ છે.
\end{enumerate}

જોકે $\lnot$ને ઔપચારિક રીતે સંક્ષેપ ગણીએ છીએ, કેટલીક વાર
તે મૂળભૂત હોય તેમ $\lnot$ માટે સ્પષ્ટ નિયમો અને વ્યાખ્યાની
શરતો આપીશું. મુખ્યત્વે અભ્યાસ-પ્રશ્નો રજૂ કરી શકાય તે માટે
આવું કરીશું.
% END OLP-0494
\endgroup


\begingroup
\makeatletter\def\ol@sectioncs{section}\makeatother

% BEGIN OLP-0495
\olfileid{int}{int}{bhk}

\olsection{બ્રાઉર--હાઇટિંગ--કોલ્મોગોરોવ અર્થઘટન}
\phantomsection\label{guunit:OLP-0495}


\begin{editorial}
  BHK અર્થઘટન વડે સ્ફુરણાવાદી વિધાનોની પ્રામાણ્યતાની
  સાબિતીઓ ગૂંચવણભરી છે; તેમને વધુ સારી રીતે સમજાવવી જોઈએ.
\end{editorial}

સ્ફુરણાવાદી સંયોજકોનું એક અનૌપચારિક રચનાત્મક અર્થઘટન છે, જે
સામાન્ય રીતે બ્રાઉર--હાઇટિંગ--કોલ્મોગોરોવ (BHK) અર્થઘટન તરીકે
ઓળખાય છે. તેમાં ``રચના''નો ખ્યાલ વપરાય છે; તેને રચનાત્મક
સાબિતી સમજી શકાય. (વિધિવત્ નિષ્પત્તિ તંત્રની
નિષ્પત્તિ સાથે ગૂંચવણ ન થાય તે માટે BHK અર્થઘટનમાં
``સાબિતી'' શબ્દ વાપરતા નથી.) આ સહજ ખ્યાલને આધારે BHK
અર્થઘટન સ્ફુરણાવાદી સંયોજકોનો અર્થ સમજાવે છે.

\begin{enumerate}
\item પરમાણુક વિધાનની રચના શેને કહેવાય તે આપણે જાણીએ છીએ એવું
  ધારી લઈએ છીએ.
\item $!A_1 \land !A_2$ની રચના એ જોડું $\tuple{M_1, M_2}$ છે,
  જેમાં $M_1$ એ~$!A_1$ની અને $M_2$ એ~$!A_2$ની રચના છે.
\item $!A_1 \lor !A_2$ની રચના એ જોડું $\tuple{s, M}$ છે:
  ક્યાં તો $s$ એ~$1$ છે અને $M$ એ~$!A_1$ની રચના છે, અથવા
  $s$ એ~$2$ છે અને $M$ એ~$!A_2$ની રચના છે.
\item $!A \lif !B$ની રચના એવું વિધેય છે જે~$!A$ની રચનાને
  $!B$ની રચનામાં ફેરવે છે.
\item $\lfalse$ (અસંગતિ)ની કોઈ રચના નથી.
\item $\lnot !A$ને $!A \lif \lfalse$ના સમાનાર્થી તરીકે
  વ્યાખ્યાયિત કરીએ છીએ. એટલે $\lnot !A$ની રચના એવું વિધેય છે
  જે~$!A$ની રચનાને~$\lfalse$ની રચનામાં ફેરવે છે.
\end{enumerate}

\begin{ex}
દાખલા તરીકે $\lnot \lfalse$ લો. તેની રચના એવું વિધેય છે જે
$\lfalse$ની કોઈ પણ રચનાને ઇનપુટ તરીકે લઈને આઉટપુટમાં
ફરી~$\lfalse$ની રચના આપે. સ્પષ્ટપણે, તાદાત્મ્ય
વિધેય~$\Id{}$ આવી રચના છે: $\lfalse$ની રચના~$M$ આપીએ,
તો $\Id{}(M) = M$ ફરી~$\lfalse$ની રચના આપે છે.
\end{ex}

સામાન્ય રીતે, $\lnot !A$નો અર્થ ``$!A$ની રચના અશક્ય છે'' એવો થાય છે.

\begin{ex}
કોઈ પણ વિધાન~$!A$ માટે $!A \lif \lnot \lnot !A$ સાબિત કરીએ;
એ $!A \lif ((!A \lif \lfalse) \lif \lfalse)$ છે. જોઈએ
એવી રચના વિધેય~$f$ છે, જે $!A$ની રચના~$M$ મળતાં
$(!A \lif \lfalse) \lif \lfalse$ની રચના~$f(M)$ આપે.
હવે $f$ એ $(!A \lif \lfalse) \lif \lfalse$ની રચના કેવી
રીતે બનાવે? આપણે એવું વિધેય $g$ વ્યાખ્યાયિત કરવું જોઈએ,
જે $!A \lif \lfalse$ની રચના~$h$ ઇનપુટ તરીકે લે અને
આઉટપુટમાં~$\lfalse$ની રચના આપે. $g$ને આ રીતે વ્યાખ્યાયિત
કરીએ: પહેલાં મળેલી $!A$ની રચના~$M$ પર ઇનપુટ~$h$ લાગુ કરો.
$h$નું આઉટપુટ~$h(M)$ એ $\lfalse$ની રચના છે; તેથી $M$
એ~$!A$ની રચના હોય તો $f(M)(h) = h(M)$ એ~$\lfalse$ની
રચના છે.
\end{ex}

\begin{ex}
$\lnot(!A \land \lnot !A)$, એટલે કે $(!A
\land (!A \lif \lfalse)) \lif \lfalse$, માટે રચના આપીએ.
આ એવું વિધેય~$f$ છે જે $!A \land (!A \lif \lfalse)$ની
રચના~$M$ ઇનપુટ તરીકે લઈને~$\lfalse$ની રચના આપે છે.
સંયોજન $!B_1 \land !B_2$ની રચના એ જોડું $\tuple{N_1, N_2}$
છે, જેમાં $N_1$ એ~$!B_1$ની અને $N_2$ એ~$!B_2$ની રચના છે.
$!B_1 \land !B_2$ની રચનામાંથી અનુક્રમે $!B_1$ અને $!B_2$ની
રચનાઓ પાછી મેળવતાં વિધેયો $p_1$ અને $p_2$ આ રીતે વ્યાખ્યાયિત
કરી શકીએ:
\begin{align*}
  p_1(\tuple{N_1, N_2}) &= N_1\\
  p_2(\tuple{N_1, N_2}) &= N_2
\end{align*}
$f$ આમ કાર્ય કરે છે: પહેલાં તે પોતાના ઇનપુટ~$M$ પર $p_1$
લાગુ કરે છે અને~$!A$ની રચના મેળવે છે. પછી $M$ પર $p_2$
લાગુ કરીને~$!A \lif \lfalse$ની રચના મેળવે છે. આ રચના
પોતે જ વિધેય~$p_2(M)$ છે, જે~$!A$ની રચનાને ઇનપુટ તરીકે
લે તો~$\lfalse$ની રચના આપે છે. બીજા શબ્દોમાં, $p_2(M)$ને
$p_1(M)$ પર લાગુ કરતાં~$\lfalse$ની રચના મળે છે. તેથી
$f(M) = p_2(M)(p_1(M))$ એમ વ્યાખ્યાયિત કરી શકીએ.
\end{ex}

\begin{ex}
$((!A \land !B) \lif !C) \lif (!A \lif (!B \lif !C))$ની
રચના આપીએ. એટલે, $(!A \land !B) \lif !C$ની રચના~$g$ને
$(!A \lif (!B \lif !C))$ની રચનામાં ફેરવતું વિધેય~$f$
આપીએ. રચના~$g$ પોતે $!A \land !B$ની રચનાઓમાંથી~$!C$ની
રચનાઓ આપતું વિધેય છે. \sourcecorrection{OLMD-083}{મૂળમાં
અહીં આઉટપુટ ``$C$ની રચનાઓ'' લખ્યું છે; આસપાસનાં બધાં
સૂત્રોમાં મેટાચલ $!C$ છે, તેથી તે જ રાખ્યું છે.}
$f(g)$નું આઉટપુટ~$h_g$ એવું વિધેય છે જે $!A$ની રચનાઓમાંથી
$!B$ની રચનાઓને~$!C$ની રચનાઓમાં ફેરવતાં વિધેયો આપે છે.

હા, આ ગૂંચવણભર્યું છે. આપણે વિધેય~$h_g$ બનાવવું છે, જે
ઇનપુટ~$g$ માટે $f$નું આઉટપુટ હશે. $h_g$નું ઇનપુટ એ~$!A$ની
રચના~$M$ છે. $h_g(M)$નું આઉટપુટ એવું વિધેય~$k_M$ હોવું જોઈએ
જે~$!B$ની રચના~$N$ને~$!C$ની રચનામાં ફેરવે. હવે
$k_{g,M}(N) = g(\tuple{M, N})$ રાખો. યાદ રાખો કે
$\tuple{M, N}$ એ~$!A \land !B$ની રચના છે. તેથી
$k_{g,M}$ એ $!B \lif !C$ની રચના છે: તે~$!B$ની
રચનાઓ~$N$ને~$!C$ની રચનાઓમાં ફેરવે છે. હવે
$h_g(M) = k_{g,M}$ રાખો. આ વિધેય~$!A$ની રચનાઓ~$M$ને
$!B \lif !C$ની રચનાઓ~$k_{g,M}$માં ફેરવે છે. અંતે
$f(g) = h_g$ રાખો. આ વિધેય $(!A \land !B) \lif !C$ની
રચનાઓ~$g$ને $!A \lif (!B \lif !C)$ની રચનાઓમાં ફેરવે
છે. હાશ!{}
\end{ex}

$!A \lor \lnot !A$ વિધાનને વર્જિત મધ્યનો નિયમ કહે છે.
ચોક્કસ~$!A$ માટે તેને સાબિત કરી શકાય છે (દાખલા તરીકે,
$\lfalse \lor \lnot \lfalse$), પરંતુ સામાન્ય રીતે નહીં.
કારણ કે સ્ફુરણાવાદી વિયોજન માટે બંનેમાંથી કોઈ એક વિયોજ્યની
રચના જરૂરી છે, જ્યારે કેટલાંક વિધાનો હાલ સાબિત પણ કરી શકાતાં
નથી અને ખંડિત પણ કરી શકાતાં નથી (જેમ કે ગોલ્ડબાખની પરિકલ્પના).
તેમ છતાં વર્જિત મધ્યના નિયમને ખંડિત પણ કરી શકાતો નથી:
એટલે $\lnot \lnot (!A \lor \lnot !A)$ સત્ય છે.

\begin{ex}
$\lnot \lnot (!A \lor \lnot !A)$ સાબિત કરવા એવું વિધેય~$f$
જોઈએ જે $\lnot(!A \lor \lnot !A)$ની, એટલે
$(!A \lor (!A \lif \lfalse)) \lif \lfalse$ની રચનાને
~$\lfalse$ની રચનામાં ફેરવે. બીજા શબ્દોમાં, એવું વિધેય~$f$
જોઈએ કે $g$ એ $\lnot(!A \lor \lnot !A)$ની રચના હોય તો
$f(g)$ એ~$\lfalse$ની રચના હોય.

ધારો કે $g$ એ $\lnot(!A \lor \lnot !A)$ની રચના છે; એટલે
એવું વિધેય છે જે $!A \lor \lnot !A$ની રચનાને~$\lfalse$ની
રચનામાં ફેરવે છે. $!A \lor \lnot !A$ની રચના એ જોડું
$\tuple{s, M}$ છે, જેમાં ક્યાં તો $s = 1$ અને $M$
એ~$!A$ની રચના છે, અથવા $s = 2$ અને $M$ એ $\lnot !A$ની
રચના છે. હવે $h_1$ એવું વિધેય હોય જે $!A$ની રચના~$M_1$ને
$!A \lor \lnot !A$ની રચનામાં ફેરવે: તે $M_1$ને
$\tuple{1, M_1}$માં મોકલે છે.
\sourcecorrection{OLMD-084}{મૂળમાં $h_1$નું ઇનપુટ $M_1$
હોવા છતાં આઉટપુટ $\tuple{1, M_2}$ લખાયું છે. વિયોજનની
પ્રથમ શાખાની રચના માટે $\tuple{1, M_1}$ મૂક્યું છે.}
એ જ રીતે $h_2$ એ $\lnot !A$ની રચના~$M_2$ને
$!A \lor \lnot !A$ની રચનામાં ફેરવતું વિધેય હોય: તે
$M_2$ને $\tuple{2, M_2}$માં મોકલે છે.

હવે $k$ને $\comp{h_1}{g}$ માનો: આ વિધેય~$!A$ની રચના
મળે તો~$\lfalse$ની રચના આપે છે; એટલે તે $!A \lif \lfalse$
અથવા $\lnot !A$ની રચના છે. હવે $l$ને $\comp{h_2}{g}$
માનો. આ વિધેયને $\lnot !A$ની રચના આપીએ તો તે~$\lfalse$ની
રચના આપે છે. $k$ એ $\lnot !A$ની રચના હોવાથી, $l(k)$
એ~$\lfalse$ની રચના છે.

આ રીતે $\lnot(!A \lor \lnot !A)$ની રચના~$g$માંથી
~$\lfalse$ની રચના કેવી રીતે મેળવવી તે બતાવ્યું છે.
એટલે $\lnot(!A \lor \lnot !A)$ની રચના~$g$ને
~$\lfalse$ની રચના~$l(k)$માં ફેરવતું વિધેય~$f$ એ
$\lnot\lnot(!A \lor \lnot !A)$ની રચના છે.
\end{ex}

જોઈ શકાય છે કે સૂત્રોની સ્ફુરણાવાદી પ્રામાણ્યતા
બતાવવા BHK અર્થઘટન વાપરવું ઝડપથી ભારરૂપ અને ગૂંચવણભર્યું
બની જાય છે. સદભાગ્યે સ્ફુરણાવાદી તર્ક માટે વધુ સારી
નિષ્પત્તિ તંત્રો અને વધુ ચોક્કસ અર્થવિજ્ઞાનલક્ષી
અર્થઘટનો ઉપલબ્ધ છે.
% END OLP-0495
\endgroup


\begingroup
\makeatletter\def\ol@sectioncs{section}\makeatother

% BEGIN OLP-0496
\olfileid{int}{int}{ntd}

\olsection{સ્વાભાવિક નિગમન}
\phantomsection\label{guunit:OLP-0496}


$\FalseCl$ નિયમો વિનાનું સ્વાભાવિક નિગમન સ્ફુરણાવાદી તર્ક
માટેનું પ્રમાણભૂત નિષ્પત્તિ તંત્ર છે. અહીં નિયમો ફરી
આપીએ છીએ અને BHK અર્થઘટન વડે તેમની પ્રેરણા સમજાવીએ છીએ.
દરેક કિસ્સામાં નિયમ એવો છે કે આધારવાક્યોની રચનાઓ હોય, તો
નિષ્કર્ષની પણ રચના હોય એવું કહી શકાય.

સ્વાભાવિક નિગમનની નિષ્પત્તિઓમાં નિવૃત્ત ન કરેલી ધારણાઓ
હોતી હોવાથી, એવી નિષ્પત્તિને---દાખલા તરીકે
અનિવૃત્ત ધારણાઓ~$\Gamma$માંથી $!A$ની નિષ્પત્તિને---
એવા વિધેય તરીકે જોવી જોઈએ જે દરેક $!B \in \Gamma$ની
રચનાઓને~$!A$ની રચનામાં ફેરવે છે. અનિવૃત્ત ધારણાઓ
વિના~$!A$ની નિષ્પત્તિ હોય, તો BHK અર્થઘટનના અર્થમાં
~$!A$ની રચના હોય છે. જોકે ચર્ચામાં જરૂર ન હોય ત્યારે
~$\Gamma$નો ઉલ્લેખ છોડીશું.

ધારણા~$!A$ પોતે જ અનિવૃત્ત ધારણા~$!A$માંથી~$!A$ની
નિષ્પત્તિ છે. આ BHK અર્થઘટનને અનુરૂપ છે: રચનાઓ પરનું
તાદાત્મ્ય વિધેય~$!A$ની કોઈ પણ રચનાને~$!A$ની રચનામાં
ફેરવે છે.

\subsection{સંયોજન}

\begin{defish}
\AxiomC{$!A$}
\AxiomC{$!B$}
\RightLabel{\Intro{\land}}
\BinaryInfC{$!A \land !B$}
\DisplayProof
\hfill
\begin{tabular}{l}
\AxiomC{$!A \land !B$}
\RightLabel{\Elim{\land}}
\UnaryInfC{$!A$}
\DisplayProof
\\[3ex]
\AxiomC{$!A \land !B$}
\RightLabel{\Elim{\land}}
\UnaryInfC{$!B$}
\DisplayProof
\end{tabular}
\end{defish}

ધારો કે $!A_1$ અને $!A_2$ની રચનાઓ અનુક્રમે $N_1$, $N_2$
છે. તો $!A_1 \land !A_2$ની રચના, એટલે જોડું
$\tuple{N_1, N_2}$, પણ આપણી પાસે છે.

BHK અર્થઘટન મુજબ $!A_1 \land !A_2$ની રચના એ જોડું
$\tuple{N_1, N_2}$ છે. \sourcecorrection{OLMD-085}{મૂળમાં
અહીં $!A_1 \land !A_1$ છે, જ્યારે આગળ $N_1$ અને $N_2$
અનુક્રમે $!A_1$ તથા $!A_2$ની રચનાઓ છે. આ સંયોજન
$!A_1 \land !A_2$ કર્યું છે.} આવું જોડું હોય, તો દરેક
સંયોજ્યની રચના પણ હોય છે: $N_1$ એ~$!A_1$ની અને $N_2$
એ $!A_2$ની રચના છે.

\subsection{શરતી સંયોજક}

\begin{defish}
  \AxiomC{$\Discharge{!A}{u}$}
  \DeduceC{$!B$}
  \DischargeRule{$\Intro{\lif}$}{u}
  \UnaryInfC{$!A \lif !B$}
  \DisplayProof
\hfill
  \AxiomC{$!A \lif !B$}
  \AxiomC{$!A$}
  \RightLabel{$\Elim{\lif}$}
  \BinaryInfC{$!B$}
  \DisplayProof
\end{defish}

અનિવૃત્ત ધારણા~$!A$માંથી~$!B$ની નિષ્પત્તિ
હોય, તો એવું વિધેય~$f$ હોય છે જે~$!A$ની રચનાઓને~$!B$ની
રચનાઓમાં ફેરવે છે. આ જ વિધેય $!A \lif !B$ની રચના છે.
તેથી $\Intro\lif$ના આધારવાક્યની રચના~$!A$ની રચના પર
આધારિત હોય, તો નિષ્કર્ષ $!A \lif !B$ની રચના હોય છે.

બીજી તરફ, ધારો કે $!A$ની રચના~$N$ અને $!A \lif !B$ની
રચના~$f$ હોય. $!A \lif !B$ની રચના એવું વિધેય છે જે
$!A$ની રચનાઓને~$!B$ની રચનાઓમાં ફેરવે છે. તેથી $f(N)$
એ~$!B$ની રચના છે; એટલે $\Elim\lif$ના નિષ્કર્ષની રચના છે.

\subsection{વિયોજન}

\begin{defish}
\begin{tabular}{l}
\AxiomC{$!A$}
\RightLabel{\Intro{\lor}}
\UnaryInfC{$!A \lor !B$}
\DisplayProof
\\[3ex]
\AxiomC{$!B$}
\RightLabel{\Intro{\lor}}
\UnaryInfC{$!A \lor !B$}
\DisplayProof
\end{tabular}
\hfill
\AxiomC{$!A \lor !B$}
\AxiomC{$\Discharge{!A}{n}$}
\DeduceC{$!C$}
\AxiomC{$\Discharge{!B}{n}$}
\DeduceC{$!C$}
\DischargeRule{\Elim{\lor}}{n}
\TrinaryInfC{$!C$}
\DisplayProof
\end{defish}

~$!A_i$ની રચના~$N_i$ હોય, તો તેમાંથી~$!A_1 \lor !A_2$ની
રચના~$\tuple{i, N_i}$ બનાવી શકીએ. બીજી તરફ, ધારો કે
~$!A_1 \lor !A_2$ની રચના છે, એટલે જોડું $\tuple{i, N_i}$
છે જેમાં $N_i$ એ~$!A_i$ની રચના છે. વળી વિધેયો $f_1$,
$f_2$ અનુક્રમે $!A_1$, $!A_2$ની રચનાઓને~$!C$ની
રચનાઓમાં ફેરવે છે. તો $f_i(N_i)$ એ~$!C$ની રચના છે,
જે~$\Elim\lor$નો નિષ્કર્ષ છે.

\subsection{અસંગતિ}

\begin{defish}
  \AxiomC{$\lfalse$}
  \RightLabel{\FalseInt}
  \UnaryInfC{$!A$}
  \DisplayProof
\end{defish}

અનિવૃત્ત ધારણાઓ~$!B_1$, \dots, $!B_n$માંથી
~$\lfalse$ની નિષ્પત્તિ હોય, તો વિધેય~$f(M_1,
\dots, M_n)$ એ $!B_1$, \dots, $!B_n$ની રચનાઓને
~$\lfalse$ની રચનામાં ફેરવે છે. $\lfalse$ની કોઈ રચના
ન હોવાથી, $!B_1$, \dots, $!B_n$ બધાની રચનાઓ પણ એકસાથે
હોઈ શકતી નથી. તેથી $f$માં એ ગુણધર્મ પણ છે કે \emph{જો}
$M_1$, \dots, $M_n$ અનુક્રમે $!B_1$, \dots, $!B_n$ની
રચનાઓ હોય, \emph{તો} $f(M_1, \dots, M_n)$ એ~$!A$ની
રચના છે.

\subsection{$\lnot$ માટેના નિયમો}

$\lnot !A$ને $!A \lif \lfalse$ તરીકે વ્યાખ્યાયિત કર્યું
હોવાથી કડક અર્થમાં~$\lnot$ માટે અલગ નિયમો જરૂરી નથી.
તેમ છતાં આપીએ તો તે નીચે જેવા હશે:

\begin{defish}
\AxiomC{$\Discharge{!A}{n}$}
\noLine
\DeduceC{$\lfalse$}
\DischargeRule{\Intro{\lnot}}{n}
\UnaryInfC{$\lnot !A$}
\DisplayProof
\hfill
\AxiomC{$\lnot !A$}
\AxiomC{$!A$}
\RightLabel{\Elim{\lnot}}
\BinaryInfC{$\lfalse$}
\DisplayProof
\end{defish}


\subsection{નિષ્પત્તિઓનાં ઉદાહરણો}

\begin{enumerate}
\item $\Proves !A \lif (\lnot !A \lif \lfalse)$, એટલે કે $\Proves !A
  \lif ((!A \lif \lfalse) \lif \lfalse)$
  \begin{prooftree}
    \AxiomC{$\Discharge{!A}{2}$}
    \AxiomC{$\Discharge{!A \lif \lfalse}{1}$}
    \RightLabel{\Elim\lif}
    \BinaryInfC{$\lfalse$}
    \DischargeRule{\Intro\lif}{1}
    \UnaryInfC{$(!A \lif \lfalse)\lif \lfalse$}
    \DischargeRule{\Intro\lif}{2}
    \UnaryInfC{$!A \lif (!A \lif \lfalse) \lif \lfalse$}
  \end{prooftree}

\item $\Proves ((!A \land !B) \lif !C) \lif (!A \lif (!B \lif !C))$
  \begin{prooftree}
    \AxiomC{$\Discharge{(!A \land !B) \lif !C}{3}$}
    \AxiomC{$\Discharge{!A}{2}$}
    \AxiomC{$\Discharge{!B}{1}$}
    \RightLabel{\Intro\land}
    \BinaryInfC{$!A \land !B$}
    \RightLabel{\Elim\lif}
    \BinaryInfC{$!C$}
    \DischargeRule{\Intro\lif}{1}
    \UnaryInfC{$!B \lif !C$}
    \DischargeRule{\Intro\lif}{2}
    \UnaryInfC{$!A \lif (!B \lif !C)$}
    \DischargeRule{\Intro\lif}{3}
    \UnaryInfC{$((!A \land !B) \lif !C) \lif (!A \lif (!B \lif !C))$}
  \end{prooftree}

\item $\Proves \lnot(!A \land \lnot !A)$, એટલે કે $\Proves (!A \land (!A
  \lif \lfalse))\lif \lfalse$
  \begin{prooftree}
    \AxiomC{$\Discharge{!A \land (!A \lif \lfalse)}{1}$}
    \RightLabel{\Elim\land}
    \UnaryInfC{$!A \lif \lfalse$}
    \AxiomC{$\Discharge{!A \land (!A \lif \lfalse)}{1}$}
    \RightLabel{\Elim\land}
    \UnaryInfC{$!A$}
    \RightLabel{\Elim\lif}
    \BinaryInfC{$\lfalse$}
    \DischargeRule{\Intro\lif}{1}
    \UnaryInfC{$(!A \land (!A \lif \lfalse))\lif \lfalse$}
  \end{prooftree}

\item $\Proves \lnot\lnot(!A \lor \lnot !A)$, એટલે કે $\Proves ((!A \lor
  (!A \lif \lfalse))\lif \lfalse)\lif \lfalse$
  \begin{prooftree}\hskip-3em
    \AxiomC{$\Discharge{(!A \lor (!A \lif \lfalse))\lif \lfalse}{2}$}
    \AxiomC{$\Discharge{(!A \lor (!A \lif \lfalse))\lif \lfalse}{2}$}
    \AxiomC{$\Discharge{!A}{1}$}
    \RightLabel{\Intro\lor}
    \UnaryInfC{$!A \lor (!A \lif \lfalse)$}
    \RightLabel{\Elim\lif}
    \BinaryInfC{$\lfalse$}
    \DischargeRule{\Intro\lif}{1}
    \UnaryInfC{$!A \lif \lfalse$}
    \RightLabel{\Intro\lor}
    \UnaryInfC{$!A \lor (!A \lif \lfalse)$}
    \RightLabel{\Elim\lif}
    \insertBetweenHyps{\hskip -4em}
    \BinaryInfC{$\lfalse$}
    \DischargeRule{\Intro\lif}{2}
    \UnaryInfC{$((!A \lor (!A \lif \lfalse))\lif \lfalse)\lif \lfalse$}
  \end{prooftree}
\end{enumerate}

\begin{prop}
  સ્ફુરણાવાદી તર્કના સ્વાભાવિક નિગમનમાં $\Gamma \Proves !A$
  હોય, તો પ્રશિષ્ટ તર્કમાં પણ $\Gamma \Proves !A$. ખાસ કરીને,
  $!A$ સ્ફુરણાવાદી પ્રમેય હોય તો પ્રશિષ્ટ પ્રમેય પણ છે.
\end{prop}

\begin{proof}
  સ્વાભાવિક નિગમનનો દરેક નિયમ પ્રશિષ્ટ સ્વાભાવિક નિગમનમાં
  પણ નિયમ છે. તેથી સ્ફુરણાવાદી તર્કની દરેક નિષ્પત્તિ
  પ્રશિષ્ટ તર્કમાં પણ નિષ્પત્તિ છે.
\end{proof}

\begin{prob}
  નીચેનાં સૂત્રોની સ્ફુરણાવાદી તર્કમાં
  નિષ્પત્તિઓ આપો:
  \begin{enumerate}
    \item $(\lnot !A \lor !B) \lif (!A \lif !B)$
    \item $\lnot\lnot\lnot !A \lif \lnot !A$
    \item $\lnot\lnot (!A \land !B) \liff (\lnot\lnot !A\land \lnot \lnot !B)$
    \item $\lnot (!A \lor !B) \liff (\lnot !A \land \lnot !B)$
    \item $(\lnot !A \lor \lnot !B) \lif \lnot (!A \land !B)$
    \item $\lnot\lnot ( !A \land !B) \lif  (\lnot\lnot !A \lor
    \lnot\lnot !B)$
  \end{enumerate}
\end{prob}
% END OLP-0496
\endgroup


\begingroup
\makeatletter\def\ol@sectioncs{section}\makeatother

% BEGIN OLP-0497
\olfileid{int}{int}{axd}

\olsection{સ્વયંસિદ્ધિમૂલક નિષ્પત્તિઓ}
\phantomsection\label{guunit:OLP-0497}


સ્ફુરણાવાદી વિધાનાત્મક તર્ક માટેની સ્વયંસિદ્ધિમૂલક
નિષ્પત્તિઓ ખ્યાલની દૃષ્ટિએ સૌથી સરળ અને ઇતિહાસમાં
સૌથી પહેલાં આવેલાં નિષ્પત્તિ તંત્રો છે. તેઓ પ્રશિષ્ટ
વિધાનાત્મક તર્કની જેમ જ કાર્ય કરે છે.

\begin{defn}[નિષ્પન્નક્ષમતા]
જો $\Gamma$ એ $\Lang L$નાં સૂત્રોનો ગણ હોય, તો
$\Gamma$માંથી \emph{નિષ્પત્તિ} એ સૂત્રોનો
સાંત ક્રમ $!A_1$, \dots,~$!A_n$ છે, જેમાં દરેક $i \le n$
માટે નીચેનામાંથી કોઈ એક શરત પૂરી થાય છે:
\begin{enumerate}
\item $!A_i \in \Gamma$; અથવા
\item $!A_i$ સ્વયંસિદ્ધિ છે; અથવા
\item $j < i$ અને $k < i$ ધરાવતા કોઈ $!A_j$ અને $!A_k$
  પરથી મોડસ પોનેન્સ વડે $!A_i$ મળે છે, એટલે કે
  $!A_k \ident !A_j \lif !A_i$.
\end{enumerate}
\end{defn}

\begin{defn}[સ્વયંસિદ્ધિઓ]
સ્ફુરણાવાદી વિધાનાત્મક તર્ક માટેની \emph{સ્વયંસિદ્ધિઓ}નો ગણ
$\PAx$ નીચેના સ્વરૂપોનાં બધાં સૂત્રોનો બનેલો છે:
\begin{align}
  & (!A \land !B) \lif !A \ollabel{ax:land1}\\
  & (!A \land !B) \lif !B \ollabel{ax:land2}\\
  & !A \lif (!B \lif (!A \land !B)) \ollabel{ax:land3}\\
  & !A \lif (!A \lor !B) \ollabel{ax:lor1}\\
  & !A \lif (!B \lor !A) \ollabel{ax:lor2}\\
  & (!A \lif !C) \lif ((!B \lif !C) \lif ((!A \lor !B) \lif !C)) \ollabel{ax:lor3}\\
  & !A \lif (!B \lif !A) \ollabel{ax:lif1}\\
  & (!A \lif (!B \lif !C)) \lif ((!A \lif !B) \lif (!A \lif !C)) \ollabel{ax:lif2}\\
  & \lfalse \lif !A \ollabel{ax:lfalse1}
\end{align}
\end{defn}

\begin{defn}[નિષ્પન્નક્ષમતા]
જો $\Gamma$માંથી અંતે $!A$ આવતું નિષ્પત્તિ હોય, તો
સૂત્ર~$!A$ એ $\Gamma$માંથી \emph{નિષ્પન્ન કરી શકાય એવું}
છે; તેને $\Gamma \Proves !A$ લખીએ છીએ.
\end{defn}

\begin{defn}[પ્રમેયો]
જો ખાલી ગણમાંથી~$!A$નું નિષ્પત્તિ હોય, તો
સૂત્ર~$!A$ એ \emph{પ્રમેય} છે. $!A$ પ્રમેય
હોય તો $\Proves !A$ અને ન હોય તો $\Proves/ !A$ લખીએ છીએ.
\end{defn}

\begin{prop}
  સ્ફુરણાવાદી તર્કના સ્વયંસિદ્ધિમૂલક નિષ્પત્તિ તંત્રમાં
  $\Gamma \Proves !A$ હોય, તો પ્રશિષ્ટ તર્કમાં પણ
  $\Gamma \Proves !A$. ખાસ કરીને, $!A$ સ્ફુરણાવાદી
  પ્રમેય હોય તો પ્રશિષ્ટ પ્રમેય પણ છે.
\end{prop}

\begin{proof}
  દરેક સ્ફુરણાવાદી સ્વયંસિદ્ધિ પ્રશિષ્ટ સ્વયંસિદ્ધિ પણ છે;
  તેથી સ્ફુરણાવાદી તર્કની દરેક નિષ્પત્તિ પ્રશિષ્ટ
  તર્કમાં પણ નિષ્પત્તિ છે.
\end{proof}
% END OLP-0497
\endgroup


\OLEndChapterHook
% END OLP-0492
\endgroup


\begingroup
\makeatletter\def\ol@sectioncs{section}\makeatother

% BEGIN OLP-0498
\olchapter{int}{sem}{અર્થવિજ્ઞાન}
\olchapterid{int}{sem}
\phantomsection\label{guunit:OLP-0498}


\begin{editorial}
  આ પ્રકરણ સ્ફુરણાવાદી તર્કના અર્થવિજ્ઞાન માટેની વ્યાખ્યાઓ
  એકત્ર કરે છે. અત્યારે તેમાં માત્ર ક્રિપ્કે અને સ્થલરચનાત્મક
  અર્થવિજ્ઞાન આવરી લેવાયાં છે. નિદર્શો સૂત્રોને કેવી રીતે સત્ય
  બનાવે છે અથવા સૂત્રોની પ્રામાણ્યતા કેવી રીતે સાબિત થાય છે
  તેનાં ઉદાહરણો હજી નથી.
\end{editorial}

\begingroup
\makeatletter\def\ol@sectioncs{section}\makeatother

% BEGIN OLP-0499
\olfileid{int}{sem}{int}

\olsection{પરિચય}
\phantomsection\label{guunit:OLP-0499}


અર્થવિજ્ઞાન વિના કોઈ તર્કપદ્ધતિનું સંતોષકારક વર્ણન થતું નથી;
સ્ફુરણાવાદી તર્ક પણ તેનો અપવાદ નથી. પ્રશિષ્ટ તર્ક માટે
સત્યમૂલ્ય-નિયુક્તિઓ પર આધારિત અર્થવિજ્ઞાન પ્રમાણભૂત છે, પરંતુ
સ્ફુરણાવાદી તર્ક માટે એકથી વધુ સ્પર્ધી અર્થવિજ્ઞાનો છે.
તેમાંનું કોઈ પણ સંયોજકોના અર્થનું સ્ફુરણાવાદી દૃષ્ટિએ
સંપૂર્ણપણે સ્વીકાર્ય વર્ણન આપે એ અર્થમાં પૂરતું સંતોષકારક નથી.

મોડલ તર્કના અર્થવિજ્ઞાન જેવું, સંબંધી નિદર્શો પર
આધારિત અર્થવિજ્ઞાન કદાચ સૌથી વધુ પ્રચલિત છે. તેમાં
પ્રતિજ્ઞાપન ચલો જગતોને સોંપવામાં આવે છે અને
આ જગતોને સુલભતા સંબંધ જોડે છે. આ સંબંધ હંમેશા આંશિક
ક્રમ હોય છે, એટલે સ્વવાચક, વિસંમિત અને પરંપરિત હોય છે.

સહજ રીતે આ જગતોને જ્ઞાનની અવસ્થાઓ અથવા ``પુરાવાની
પરિસ્થિતિઓ'' માની શકાય. અવસ્થા~$w'$ એ~$w$માંથી ત્યારે અને
માત્ર ત્યારે જ સુલભ છે જ્યારે આપણને જાણીતું છે તે મુજબ
$w'$ એ જ્ઞાનની શક્ય (ભાવિ) અવસ્થા હોય, એટલે~$w$માં
જે જાણીતું છે તેની સાથે સુસંગત હોય. એક વાર કોઈ વિધાનની
ખબર પડે, તો તે અજાણ્યું બની શકતું નથી: $!A$ એ~$w$માં
જાણીતું હોય અને~$Rww'$ હોય, તો $!A$ એ~$w'$માં પણ
જાણીતું હોય. તેથી સુલભતા સંબંધની દૃષ્ટિએ ``જ્ઞાન''
એકદિશવર્ધી છે.

જ્ઞાનસંબંધી તર્કની જેમ ``$!A$ જાણીતું છે''નો અર્થ
``બધા જ્ઞાનસંબંધી વિકલ્પોમાં સત્ય છે'' એવો લઈએ, તો
બધા વિકલ્પોમાં $!A$ અને~$!B$ બંને જાણીતાં હોય ત્યારે
$!A \land !B$ એ~$w$માં જાણીતું છે. પરંતુ જ્ઞાન
એકદિશવર્ધી છે અને $R$ સ્વવાચક છે; તેથી $!A \land !B$
એ~$w$માં ત્યારે અને માત્ર ત્યારે જ જાણીતું હોય છે જ્યારે
$!A$ અને $!B$ બંને $w$માં જાણીતાં હોય. એ જ કારણસર
$!A \lor !B$ એ~$w$માં ત્યારે અને માત્ર ત્યારે જ
જાણીતું હોય છે જ્યારે બેમાંથી ઓછામાં ઓછું એક જાણીતું હોય.
તેથી $\land$ અને $\lor$ માટે સંયોજકોની સત્ય-શરતો
પ્રશિષ્ટ તર્ક જેવી જ છે.

પરંતુ શરતી સંયોજકની સત્ય-શરતો પ્રશિષ્ટ તર્કથી જુદી છે.
$!A \lif !B$ એ~$w$માં ત્યારે અને માત્ર ત્યારે જ જાણીતું
હોય છે જ્યારે $Rww'$ ધરાવતી કોઈ પણ અવસ્થા~$w'$માં $!A$
જાણીતું હોય અને સાથે $!B$ જાણીતું ન હોય એવું ન બને.
આ શરત~$w$માં $!A$ અજાણ્યું હોય અથવા $!B$ જાણીતું
હોય તે શરત જેવી નથી. કારણ કે~$w$માં $!A$ કે $!B$
બેમાંથી કોઈ જાણીતું ન હોય, છતાં $Rww'$ ધરાવતી ભાવિ
જ્ઞાનસંબંધી અવસ્થા~$w'$ હોઈ શકે; એ $w'$માં $!A$
જાણીતું થાય અને $!B$ જાણીતું ન થાય એવું બની શકે છે.

$\lnot !A$ની ખબર આપણને ત્યારે જ હોય જ્યારે કોઈ શક્ય
ભાવિ જ્ઞાનસંબંધી અવસ્થામાં~$!A$ની ખબર પડી શકે તેમ ન હોય.
વિચાર એ છે કે $!A$ જાણવું શક્ય હોય, તો કોઈ શક્ય ભાવિ
અવસ્થામાં~$!A$ જાણીતું બને. $\lfalse$ની ખબર પડી શકતી
નથી, એટલે એવી ભાવિ અવસ્થામાં~$!A$ની ખબર હોય પણ
~$\lfalse$ની ન હોય.

આ અર્થઘટનમાં વર્જિત મધ્યનો નિયમ નિષ્ફળ જાય છે. કેટલાંક
$!A$ અત્યારે આપણને જાણીતાં નથી, પરંતુ આગળ જતાં જાણી
શકીએ. એવા સૂત્ર~$!A$ માટે $!A$ અને
~$\lnot !A$ બંને અજાણ્યાં છે, તેથી $!A \lor \lnot !A$
પણ જાણીતું નથી. પરંતુ, દાખલા તરીકે, $\lnot(!A \land
\lnot !A)$ જાણીતું છે: $!A$ અને~$\lnot !A$ બંનેની
ખબર હોય એવી કોઈ ભાવિ અવસ્થા શક્ય નથી. આપણને હાલ $!A$
કે $\lnot !A$ની ખબર છે કે નહીં તેનાથી સ્વતંત્ર રીતે આ
વાત જાણીએ છીએ.

સંબંધી નિદર્શો જ સ્ફુરણાવાદી તર્ક માટે ઉપલબ્ધ
એકમાત્ર અર્થવિજ્ઞાન નથી. સ્થલરચનાત્મક અર્થવિજ્ઞાન બીજું
છે: તેમાં વિધાનોને સ્થલરચનાત્મક અવકાશમાં ખુલ્લા ગણો
તરીકે અર્થઘટિત કરીએ છીએ, અને સંયોજકોને આ ગણો પરની
સંક્રિયાઓ તરીકે (દાખલા તરીકે, $\land$ને છેદને અનુરૂપ)
અર્થઘટિત કરીએ છીએ.
% END OLP-0499
\endgroup

\begingroup
\makeatletter\def\ol@sectioncs{section}\makeatother

% BEGIN OLP-0500
\olfileid{int}{sem}{rel}

\olsection{સંબંધી નિદર્શો}
\phantomsection\label{guunit:OLP-0500}


સ્ફુરણાવાદી વિધાનાત્મક તર્કનું ચોક્કસ અર્થવિજ્ઞાન આપવા,
જેના સંદર્ભમાં સૂત્રોનું મૂલ્યાંકન કરી શકીએ એવા
નિદર્શની વ્યાખ્યા આપવી પડે. એવી વ્યાખ્યાને આધારે પછી
પ્રામાણ્યતા અને અનુસરણ જેવા અર્થવિજ્ઞાનલક્ષી ખ્યાલો પણ
વ્યાખ્યાયિત કરી શકાય છે. આવું એક અર્થવિજ્ઞાન
સંબંધી નિદર્શો વડે મળે છે.

\begin{defn}
  સ્ફુરણાવાદી વિધાનાત્મક તર્કનું સંબંધી નિદર્શ
  એ ત્રિક $\mModel{M} = \tuple{W, R, V}$ છે, જ્યાં
  \begin{enumerate}
  \item $W$ એ અખાલી ગણ છે,
  \item $R$ એ~$W$ પરનો આંશિક ક્રમ (એટલે સ્વવાચક, વિસંમિત
    અને પરંપરિત દ્વિસ્થાની સંબંધ) છે, અને
  \item $V$ એવું વિધેય છે જે દરેક પ્રતિજ્ઞાપન ચલ~$p$ને~$W$નો ઉપગણ આપે છે, જેથી
  \item $V$ એ $R$ની દૃષ્ટિએ એકદિશવર્ધી છે: જો $w
    \in V(p)$ અને $Rww'$ હોય, તો $w' \in V(p)$.
  \end{enumerate}
\end{defn}

\begin{defn}\ollabel{defn:true-at-w}
  નિદર્શ~$\mModel{M}$માં $!A$નું \emph{$w$ ખાતે સત્ય
    હોવું}, સંજ્ઞામાં $\mSat{M}{!A}[w]$, નીચે મુજબ
  આગમનાત્મક રીતે વ્યાખ્યાયિત કરીએ છીએ:
  \begin{enumerate}
  \item \indcase{!A}{p}{$\mSat{M}{\indfrm}[w]$ ત્યારે અને માત્ર ત્યારે જ જ્યારે $w \in V(p)$.}
  \item \indcase{!A}{\lfalse}{$\mSat{M}{\indfrm}[w]$ નહીં}.
  \item \indcase{!A}{\lnot !B}{$\mSat{M}{\indfrm}[w]$ ત્યારે અને માત્ર ત્યારે જ જ્યારે
    $Rww'$ ધરાવતું એવું કોઈ $w'$ નથી કે $\mSat{M}{!B}[w']$}.
  \item \indcase{!A}{!B \land !C}{$\mSat{M}{\indfrm}[w]$ ત્યારે અને માત્ર ત્યારે જ જ્યારે
    $\mSat{M}{!B}[w]$ અને $\mSat{M}{!C}[w]$}.
  \item \indcase{!A}{!B \lor !C}{$\mSat{M}{\indfrm}[w]$ ત્યારે અને માત્ર ત્યારે જ જ્યારે
    $\mSat{M}{!B}[w]$ અથવા $\mSat{M}{!C}[w]$ (અથવા બંને)}.
  \item \indcase{!A}{!B \lif !C}{$\mSat{M}{\indfrm}[w]$ ત્યારે અને માત્ર ત્યારે જ જ્યારે
    $Rww'$ ધરાવતા દરેક $w'$ માટે $\mSat{M}{!B}[w']$ નહીં
    અથવા $\mSat{M}{!C}[w']$} (અથવા બંને).
  \end{enumerate}
  $\mSat{M}{!A}[w]$ ન હોય તો $\mSat/{M}{!A}[w]$ લખીએ છીએ.
  $\Gamma$ એ સૂત્રોનો ગણ હોય, તો
  $\mSat{M}{\Gamma}[w]$નો અર્થ બધા $!B \in \Gamma$
  માટે $\mSat{M}{!B}[w]$ એવો થાય છે.
\end{defn}

\begin{prob}
  \olref[int][sem][rel]{defn:true-at-w} મુજબ
  $\mSat{M}{\lnot !A}[w]$ ત્યારે અને માત્ર ત્યારે જ જ્યારે
  $\mSat{M}{!A \lif \lfalse}[w]$ થાય છે તે બતાવો.
\end{prob}

\begin{prop}\ollabel{prop:true-monotonic}
  જગતોમાં સત્યતા~$R$ની દૃષ્ટિએ એકદિશવર્ધી છે: જો
  $\mSat{M}{!A}[w]$ અને $Rww'$ હોય, તો
  $\mSat{M}{!A}[w']$.
\end{prop}

\begin{proof}
  સ્વાધ્યાય.
\end{proof}

\begin{prob}
  \olref[int][sem][rel]{prop:true-monotonic} સાબિત કરો.
\end{prob}
% END OLP-0500
\endgroup

\begingroup
\makeatletter\def\ol@sectioncs{section}\makeatother

% BEGIN OLP-0501
\olfileid{int}{sem}{sem}

\olsection{અર્થવિજ્ઞાનલક્ષી ખ્યાલો}
\phantomsection\label{guunit:OLP-0501}


\begin{defn}
  બધા $w \in W$ માટે $\mSat{M}{!A}[w]$ હોય ત્યારે અને માત્ર
  ત્યારે જ $!A$ એ \emph{નિદર્શ $\mModel{M} = \tuple{W,R,V}$માં
  સત્ય છે}, એટલે $\mSat{M}{!A}$, એમ કહીએ છીએ.
  બધા નિદર્શોમાં સત્ય હોય ત્યારે અને માત્ર ત્યારે જ $!A$
  \emph{પ્રામાણ્ય} ધરાવે છે, એટલે $\Entails !A$.
  સૂત્રોનો ગણ~$\Gamma$ એ~$!A$ને \emph{ફલિત કરે છે},
  એટલે $\Gamma \Entails !A$, એમ ત્યારે અને માત્ર ત્યારે જ
  કહીએ છીએ જ્યારે દરેક નિદર્શ~$\mModel{M}$ અને દરેક~$w$
  માટે $\mSat{M}{\Gamma}[w]$ હોય તો $\mSat{M}{!A}[w]$ હોય.
\end{defn}

\begin{prop}\ollabel{prop:sat-entails}
  \begin{enumerate}
  \item\ollabel{prop:sat-entails1} જો $\mSat{M}{\Gamma}[w]$ અને
    $\Gamma \Entails !A$ હોય, તો $\mSat{M}{!A}[w]$.
  \item\ollabel{prop:sat-entails2} જો $\mSat{M}{\Gamma}$ અને $\Gamma
    \Entails !A$ હોય, તો $\mSat{M}{!A}$.
  \end{enumerate}
\end{prop}

\begin{proof}
  \sourcecorrection{OLMD-086}{મૂળમાં પહેલી પંક્તિની સાબિતી
  બીજી પંક્તિની સર્વજગત શરત માનીને આપે છે, અને પછી બીજી
  પંક્તિને પહેલીમાંથી અનુસરતી કહે છે. પહેલી પંક્તિ વ્યાખ્યાથી
  સીધી મળે છે; મૂળનો સર્વજગત તર્ક બીજી પંક્તિમાં મૂક્યો છે.}
\begin{enumerate}

  \item પહેલું વિધાન અનુસરણની વ્યાખ્યાથી સીધું મળે છે.
  \item ધારો કે $\mSat{M}{\Gamma}$. $\Gamma \Entails !A$
    હોવાથી $\mSat{M}{\Gamma}[w]$ હોય તો
    $\mSat{M}{!A}[w]$ થાય છે. દરેક $u \in W$ માટે
    $\mSat{M}{\Gamma}[u]$ હોવાથી
    $\mSat{M}{\Gamma}[w]$ પણ છે. તેથી
    $\mSat{M}{!A}[w]$. દરેક જગત માટે આ સાચું છે;
    આ \olref{prop:sat-entails1}ને દરેક જગતમાં લાગુ કરીને
    પણ તરત મળે છે.
  \end{enumerate}
\end{proof}


\begin{defn}\ollabel{defn:restrict}
  ધારો કે $\mModel{M}$ એ સંબંધાત્મક નિદર્શ છે અને $w \in W$.
  $\mModel{M}$નું~$w$ સુધીનું \emph{મર્યાદન}
  $\mModel{M}_w=\tuple{W_w, R_w, V_w}$ આ પ્રમાણે છે:
  \begin{align*}
    W_w & = \Setabs{u \in W}{Rwu},\\
    R_w  & = R \cap (W_w)^2, \text{ અને}\\
    V_w(p) & = V(p) \cap W_w.
  \end{align*}
\end{defn}

\begin{prop}\ollabel{prop:restrict}
  $\mSat{M}{!A}[w]$ ત્યારે અને માત્ર ત્યારે જ જ્યારે
  $\mSat{M_w}{!A}$.
\end{prop}

\begin{prob}
  \olref[int][sem][sem]{prop:restrict} સાબિત કરો.
\end{prob}

\begin{prop}
  ધારો કે $\mSat{M}{\Gamma}$ ધરાવતા દરેક નિદર્શ~$\mModel{M}$
  માટે $\mSat{M}{!A}$. તો $\Gamma \Entails !A$.
\end{prop}

\begin{proof}
  ધારો કે $\mSat{M}{\Gamma}[w]$. દરેક~$!B \in \Gamma$
  પર \olref{prop:restrict} લાગુ કરતાં $\mSat{M_w}{\Gamma}$
  મળે છે. ધારણા મુજબ $\mSat{M_w}{!A}$. ફરી
  \olref{prop:restrict} વડે $\mSat{M}{!A}[w]$ મળે છે.
\end{proof}
% END OLP-0501
\endgroup

\begingroup
\makeatletter\def\ol@sectioncs{section}\makeatother

% BEGIN OLP-0502
\olfileid{int}{sem}{top}

\olsection{સ્થલરચનાત્મક અર્થવિજ્ઞાન}
\phantomsection\label{guunit:OLP-0502}


સ્ફુરણાવાદી તર્ક માટે અર્થવિજ્ઞાન આપવાની બીજી રીત
સ્થલરચનાના ગાણિતિક ખ્યાલનો ઉપયોગ કરવાની છે.

\begin{defn}
  ધારો કે $X$ એક ગણ છે. $X$ પરની \emph{સ્થલરચના} એ
  $\Top{O} \subseteq \Pow{X}$ એવો ગણ છે જે નીચેના ગુણધર્મો
  સંતોષે છે. $\Top{O}$ના ઘટકોને સ્થલરચનાના
  \emph{ખુલ્લા ગણો} કહે છે. $X$ અને $\Top{O}$ને સાથે મળીને
  \emph{સ્થલરચનાત્મક અવકાશ} કહે છે.
  \begin{enumerate}
  \item ખાલી ગણ અને આખો અવકાશ ખુલ્લા છે: $\emptyset$, $X \in
    \Top{O}$.
  \item ખુલ્લા ગણો સાન્ત છેદ હેઠળ બંધ છે: જો $U$, $V \in
    \Top{O}$ હોય, તો $U \cap V \in \Top{O}$.
  \item ખુલ્લા ગણો મનસ્વી યોગગણ હેઠળ બંધ છે: જો બધા $i \in I$
    માટે $U_i \in \Top{O}$ હોય, તો $\bigcup \Setabs{U_i}{i \in
      I} \in \Top{O}$.
  \end{enumerate}
\end{defn}

ખુલ્લા ગણોનો સંગ્રહ સંદર્ભમાંથી જાણી શકાય ત્યારે આપણે
સ્થલરચનાત્મક અવકાશ માટે માત્ર $X$ લખી શકીએ છીએ. જોકે,
$X$ પર ખુલ્લા ગણો નક્કી કર્યા પછી જ તેને સ્થલરચનાત્મક
અવકાશ કહી શકાય.
\sourcecorrection{OLMD-087}{મૂળ આ વાક્યમાં $X$ને
``સ્થલરચના'' કહે છે; ઉપરની વ્યાખ્યા મુજબ સ્થલરચના
$\Top{O}$ છે, જ્યારે $X$ અને $\Top{O}$ સાથે મળીને
સ્થલરચનાત્મક અવકાશ બને છે.}

\begin{defn}
  સ્ફુરણાવાદી વિધાનાત્મક તર્કનું \emph{સ્થલરચનાત્મક નિદર્શ}
  એ ત્રિક $\mModel{X} = \tuple{X, \Top{O}, V}$ છે, જ્યાં
  $\Top{O}$ એ~$X$ પરની સ્થલરચના છે અને $V$ એ એવું વિધેય
  છે જે દરેક વિધાનાત્મક ચલને $\Top{O}$નો એક ખુલ્લો ગણ આપે છે.

  આપેલું સ્થલરચનાત્મક નિદર્શ~$\mModel{X}$ હોય, તો
  $\Prop{X}{!A}$ને નીચે મુજબ આગમનાત્મક રીતે વ્યાખ્યાયિત કરીએ છીએ:
  \begin{enumerate}
  \item $\Prop{X}{\lfalse} = \emptyset$
  \item $\Prop{X}{p} = V(p)$
  \item $\Prop{X}{!A \land !B} = \Prop{X}{!A} \cap \Prop{X}{!B}$
  \item $\Prop{X}{!A \lor !B} = \Prop{X}{!A} \cup \Prop{X}{!B}$
  \item $\Prop{X}{!A \lif !B} = \Interior{(X \setminus \Prop{X}{!A}) \cup
    \Prop{X}{!B}}$
  \end{enumerate}
  અહીં $\Interior{V}$ એ $V \subseteq X$ ગણને તેના
  \emph{અંતર્ભાગ} પર લઈ જતું વિધેય છે; એટલે, તેમાં સમાયેલા
  બધા ખુલ્લા ગણોનો યોગગણ. બીજા શબ્દોમાં,
  \[
  \Interior{V} = \bigcup \Setabs{U}{U \subseteq V \text{ અને } U \in \Top{O}}.
  \]
\end{defn}

નોંધો કે કોઈ પણ ગણનો અંતર્ભાગ હંમેશા ખુલ્લો છે, કારણ કે
તે ખુલ્લા ગણોનો યોગગણ છે. તેથી $\Prop{X}{!A}$ હંમેશા
ખુલ્લો ગણ છે.

સ્થલરચનાત્મક અર્થવિજ્ઞાન ખૂબ અમૂર્ત છે, છતાં તેને સમજવાની
કેટલીક રીતો તેની પાછળનો વિચાર સ્પષ્ટ કરી શકે છે. ધારો કે
$X$ના ઘટકો, અથવા ``બિંદુઓ'', એવા બિંદુઓ છે જ્યાં
વિધાનોનું મૂલ્યાંકન કરી શકાય છે. $!A$ જ્યાં સાચું હોય એવા
બધા બિંદુઓનો ગણ એ~$!A$ દ્વારા વ્યક્ત થયેલું વિધાન છે.
બિંદુઓનો દરેક ગણ શક્ય વિધાન નથી; ફક્ત $\Top{O}$ના
ઘટકો જ શક્ય વિધાનો છે. દરેક~$X$ માટે $!A \Entails !B$
ત્યારે અને માત્ર ત્યારે જ જ્યારે $!A$ જ્યાં સાચું હોય તે
દરેક બિંદુએ $!B$ પણ સાચું હોય, એટલે $\Prop{X}{!A}
\subseteq \Prop{X}{!B}$. અસંગત વિધાન~$\lfalse$ ક્યારેય
સાચું નથી, તેથી $\Prop{X}{\lfalse} = \emptyset$.

$!B \land !C$, $!B \lor !C$ અને $!B \lif !C$ દ્વારા
વ્યક્ત થયેલાં વિધાનોનો, $!B$ અને~$!C$ દ્વારા વ્યક્ત થયેલાં
વિધાનો સાથે કયો સંબંધ હોવો જોઈએ જેથી સ્ફુરણાવાદી રીતે
પ્રામાણ્ય નિયમો જળવાય, એટલે $!A \Proves !B$ ત્યારે અને માત્ર
ત્યારે જ જ્યારે $\Prop{X}{!A} \subseteq \Prop{X}{!B}$?
\sourcecorrection{OLMD-088}{મૂળમાં આ સમકક્ષતા અને અંતિમ
ફકરાની શરતમાં $\subset$ છે; સમાન ગણો માટે પણ અનુસરણ
શક્ય હોવાથી બંને જગ્યાએ $\subseteq$ કર્યું છે.}
કોઈ પણ $!A$ માટે $\lfalse \Proves !A$ જોઈએ, અને બધા
$U$ માટે $\emptyset \subseteq U$ હોવાથી આ શરત સંતોષાય
છે. $!B \land !C \Proves !B$ હોવાથી
$\Prop{X}{!B \land !C} \subseteq \Prop{X}{!B}$ જોઈએ;
એ જ રીતે $\Prop{X}{!B \land !C} \subseteq \Prop{X}{!C}$
જોઈએ. $W \subseteq U$ અને $W \subseteq V$ સંતોષતો સૌથી
મોટો ગણ $U \cap V$ છે. ઊલટું, $!B \Proves !B \lor !C$
અને $!C \Proves !B \lor !C$ હોવાથી
$\Prop{X}{!B} \subseteq \Prop{X}{!B \lor !C}$ અને
$\Prop{X}{!C} \subseteq \Prop{X}{!B \lor !C}$ જોઈએ.
$U \subseteq W$ અને $V \subseteq W$ સંતોષતો સૌથી નાનો
ગણ~$W$ એ $U \cup V$ છે.

$\lif$ માટેની વ્યાખ્યા થોડી મુશ્કેલ છે: $!A \lif !B$ એ
એવું સૌથી નબળું વિધાન વ્યક્ત કરે છે જે $!A$ સાથે મળીને
$!B$ને ફલિત કરે. $!A \lif !B$ અને $!A$ સાથે મળીને~$!B$
ફલિત થાય છે તે $(!A \lif !B) \land !A \Proves !B $
પરથી સ્પષ્ટ છે. તેથી $\Prop{X}{!A \lif !B}$ એવો સૌથી મોટો
ખુલ્લો ગણ હોવો જોઈએ કે $\Prop{X}{!A \lif !B} \cap \Prop{X}{!A}
\subseteq \Prop{X}{!B}$; આથી આપણી વ્યાખ્યા મળે છે.
% END OLP-0502
\endgroup


\OLEndChapterHook


\begin{editorial}
આગળનો સંબંધિત અંશ મૂળના સક્રિય આયાતક્રમમાં નથી. તેને અહીં સ્પષ્ટ વૈકલ્પિક સંદર્ભમાં જાળવ્યો છે.
\end{editorial}
\begingroup
\makeatletter\def\ol@sectioncs{section}\makeatother

% BEGIN OLP-0649
\olfileid{int}{sem}{pro}

\olsection{વિધાનો}
\phantomsection\label{guunit:OLP-0649}


ક્યારેક સંબંધી નિદર્શ~$\mModel{M}$માં
સૂત્ર~$!A$ જ્યાં સાચું હોય તે જગતોના
ગણને નામ આપવું ઉપયોગી છે. તેને $!A$ દ્વારા
વ્યાખ્યાયિત \emph{વિધાન} કહીએ છીએ અને
$\Prop{M}{!A}$ લખીએ છીએ.

\begin{defn}
  $V$ વિધેયને આગમનાત્મક રીતે
  $\Prop{M}{\cdot} \colon \Frm
  \to \Pow{W}$ વિધેય સુધી આ પ્રમાણે વિસ્તરીએ:
  \begin{enumerate}
  \item $\Prop{M}{\lfalse} = \emptyset$
  \item $\Prop{M}{p} = V(p)$
  \item $\Prop{M}{\lnot !B} = {}$
    \[
      \Setabs{w \in W}{\text{દરેક $w'$ માટે જે $Rww'$ સંતોષે છે, $w'
          \notin \Prop{M}{!B}$}}.
    \]
  \item $\Prop{M}{!B \land !C} = \Prop{M}{!B} \cap \Prop{M}{!C}$.
  \item $\Prop{M}{!B \lor !C} = \Prop{M}{!B} \cup \Prop{M}{!C}$.
  \item $\Prop{M}{!B \lif !C} = {}$
    \[
      \Setabs{w \in W}{\text{દરેક $w'$ માટે જે $Rww'$ સંતોષે છે, જો $w'
          \in \Prop{M}{!B}$ હોય તો $w' \in \Prop{M}{!C}$}}.
    \]
  \end{enumerate}
  \sourcecorrection{OLMD-185}{મૂળમાં આખા વિધેયની
  પ્રકાર-ઘોષણામાં $\Prop{M}{!A}$ લખ્યું છે, જોકે
  $!A$ અહીં વિધેયનું આર્ગ્યુમેન્ટ છે; તેથી
  $\Prop{M}{\cdot}$ રાખ્યું છે.}
\end{defn}

\begin{prop}
  \ollabel{prop:propositions-true}
  $\mSat{M}{!A}[w]$ ત્યારે અને માત્ર ત્યારે જ જ્યારે
  $w \in \Prop{M}{!A}$.
  \sourcecorrection{OLMD-186}{મૂળ વિધાનમાં
  $\mModel{M}{!A}[w]$ છે; અગાઉ વ્યાખ્યાયિત
  જગત-ખાતે સત્યતાની સંજ્ઞા $\mSat{M}{!A}[w]$
  અહીં જરૂરી છે.}
\end{prop}

\begin{proof}
  સ્વાધ્યાય.
\end{proof}

\begin{prob}
  \olref[int][sem][pro]{prop:propositions-true} સાબિત કરો.
\end{prob}
% END OLP-0649
\endgroup
% END OLP-0498
\endgroup


\begingroup
\makeatletter\def\ol@sectioncs{section}\makeatother

% BEGIN OLP-0503
\olchapter{int}{sc}{યથાર્થતા અને પૂર્ણતા}
\olchapterid{int}{sc}
\phantomsection\label{guunit:OLP-0503}


\begin{editorial}
  આ પ્રકરણ વિધાનાત્મક સ્ફુરણાવાદી તર્ક માટેના યથાર્થતા
  અને પૂર્ણતાના પરિણામો એકત્ર કરે છે. તેમાં પરિચય ઉમેરવાની
  જરૂર છે. પૂર્ણતાની સાબિતીમાં સાબિત કરી શકાય તેવા હોવા
  વિશેનાં કેટલાંક તથ્યો વપરાય છે; તેમને ક્યાંક સ્પષ્ટ રીતે
  જણાવીને સાબિત કરવાં જોઈએ.
\end{editorial}

\begingroup
\makeatletter\def\ol@sectioncs{section}\makeatother

% BEGIN OLP-0504
\olfileid{int}{sc}{sax}

\olsection{સ્વયંસિદ્ધિમૂલક નિષ્પત્તિઓની યથાર્થતા}
\phantomsection\label{guunit:OLP-0504}


\begin{editorial}
  યથાર્થતાની સાબિતી દરેક સ્વયંસિદ્ધિ સ્ફુરણાવાદી રીતે
  પ્રામાણ્ય ધરાવે છે એ હકીકત પર આધાર રાખે છે; આ હજી
  સાબિત કરવાનું બાકી છે, દાખલા તરીકે અર્થવિજ્ઞાનના પ્રકરણમાં.
\end{editorial}

\begin{thm}[યથાર્થતા]
  \ollabel{thm:soundness}
  જો $\Gamma \Proves !A$ હોય, તો $ \Gamma \Entails !A$.
\end{thm}

\begin{proof}
  જો $\Gamma \Proves !A$ હોય, તો $\Gamma \Entails !A$
  તે સાબિત કરીએ છીએ. $\Gamma$માંથી~$!A$ના નિષ્પત્તિમાં
  આવતા સૂત્રોની સંખ્યા~$n$ પર આગમન કરીએ છીએ.
  $!A_1$, \dots, $!A_n = !A$ એ~$\Gamma$માંથી
  નિષ્પત્તિ હોય, તો $\Gamma \Entails !A_n$ છે તે બતાવીએ.
  નોંધો કે $!A_1$, \dots, $!A_n$ એ નિષ્પત્તિ હોય,
  તો કોઈ પણ $k<n$ માટે $!A_1$, \dots, $!A_k$ પણ
  નિગમન છે.

  લંબાઈ~$0$નું કોઈ નિષ્પત્તિ નથી, તેથી $n=0$ માટે
  દાવો નિર્વિકલ્પે સાચો છે. હવે લંબાઈ~$<n$નાં બધાં
  નિષ્પત્તિઓ માટે દાવો સાચો માની લઈએ. $!A_n$નું
  સમર્થન કેવી રીતે મળે છે તે મુજબ કિસ્સા અલગ પાડીએ છીએ.

  \begin{enumerate}
  \item $!A_n$ એક સ્વયંસિદ્ધિ છે. બધી સ્વયંસિદ્ધિઓ
    પ્રામાણ્ય ધરાવે છે, તેથી કોઈ પણ~$\Gamma$ માટે
    $\Gamma \Entails !A_n$.
  \item $!A_n \in \Gamma$. ત્યારે કોઈ પણ $\mModel{M}$
    અને $w$ માટે, જો $\mSat{M}{\Gamma}[w]$ હોય, તો
    સ્પષ્ટ છે કે $\mSat{M}{!A_n}[w]$; એટલે
    $\Gamma \Entails !A$.
    \sourcecorrection{OLMD-089}{મૂળમાં અહીં
    $\mSat{M}{\Gamma}{!A_n}[w]$ એમ ત્રણ દલીલવાળી
    ખામીયુક્ત સંજ્ઞા છે; જરૂરી નિષ્કર્ષ
    $\mSat{M}{!A_n}[w]$ છે.}
  \item $!A_i$ અને $!A_j \ident !A_i \lif !A_n$ પરથી
    \MP{} વડે $!A_n$ મળે છે. $!A_1$, \dots, $!A_i$
    અને $!A_1$, \dots, $!A_j$ એ~$\Gamma$માંથી
    નિષ્પત્તિઓ છે, તેથી આગમનની ધારણાથી
    $\Gamma \Entails !A_i$ અને
    $\Gamma \Entails !A_i \lif !A_n$.

    ધારો કે $\mSat{M}{\Gamma}[w]$. $\mSat{M}{\Gamma}[w]$
    અને
    $\Gamma \Entails !A_i \lif !A_n$ પરથી
    $\mSat{M}{!A_i \lif !A_n}[w]$ મળે છે. વ્યાખ્યા મુજબ,
    $Rww'$ હોય એવા દરેક $w'$ માટે, જો
    $\mSat{M}{!A_i}[w']$ હોય, તો
    $\mSat{M}{!A_n}[w']$ થાય છે. $R$ સ્વવાચક હોવાથી
    $Rww'$ ધરાવતા $w'$માં $w$ પોતે પણ આવે છે;
    એટલે જો $\mSat{M}{!A_i}[w]$ હોય, તો
    $\mSat{M}{!A_n}[w]$. હવે $\Gamma \Entails !A_i$
    હોવાથી $\mSat{M}{!A_i}[w]$ છે. તેથી ઇચ્છિત
    $\mSat{M}{!A_n}[w]$ મળે છે.
  \end{enumerate}
\end{proof}
% END OLP-0504
\endgroup

\begingroup
\makeatletter\def\ol@sectioncs{section}\makeatother

% BEGIN OLP-0505
\olfileid{int}{sc}{snd}

\olsection{સ્વાભાવિક નિગમનની યથાર્થતા}
\phantomsection\label{guunit:OLP-0505}


હવે આપણે સંબંધાત્મક અર્થવિજ્ઞાનની દૃષ્ટિએ સ્વાભાવિક નિગમનની
યથાર્થતા સાબિત કરીશું. એટલે કે, ધારણાઓના ગણમાંથી
સૂત્ર નિષ્પન્ન કરી શકાય એવું હોય, તો તે ધારણાઓનો ગણ તે
સૂત્રને ફલિત કરે છે એમ બતાવીશું.

\begin{thm}[યથાર્થતા]
  \ollabel{thm:soundness}
  જો $\Gamma \Proves !A$ હોય, તો $ \Gamma \Entails !A$.
\end{thm}

\begin{proof}
  જો $\Gamma \Proves !A$ હોય, તો $\Gamma \Entails !A$
  તે સાબિત કરીએ છીએ. $\Gamma$માંથી~$!A$ના
  નિષ્પત્તિ પર આગમન કરીએ છીએ.

  \begin{enumerate}
  \item જો નિષ્પત્તિમાં માત્ર ધારણા~$!A$ જ હોય,
    તો $!A \Proves !A$ છે અને $!A \Entails !A$
    બતાવવાનું છે. ધારો કે $\mSat{M}{!A}[w]$.
    તો સ્પષ્ટ છે કે $\mSat{M}{!A}[w]$.

  \item નિષ્પત્તિનો છેલ્લો નિયમ $\Intro{\land}$ છે:
    અનિવૃત્ત ધારણાઓ~$\Gamma$માંથી
      $!B$ની અને અનિવૃત્ત ધારણાઓ~$\Delta$માંથી~$!C$ની
      પૂર્વધારણાઓના નિષ્પત્તિઓ,
      $\Gamma \Proves !B$ અને $\Delta \Proves !C$ બતાવે છે.
      આગમનની ધારણાથી $\Gamma \Entails !B$ અને
      $\Delta \Entails !C$ છે. આખા નિષ્પત્તિની
      અનિવૃત્ત ધારણાઓ $\Gamma$ અને $\Delta$ બંને
      છે, તેથી $\Gamma \cup \Delta \Entails !B \land !C$
      બતાવવાનું છે.
      \sourcecorrection{OLMD-090}{મૂળમાં અહીં નિષ્કર્ષ
      $!A \land !B$ લખાયો છે, જ્યારે પૂર્વધારણાઓ
      $!B$ અને $!C$ છે અને અંતે $!B \land !C$ મળે છે.}
      ધારો કે $\mSat{M}{\Gamma \cup \Delta}[w]$.
      ત્યારે $\mSat{M}{\Gamma}[w]$ પણ છે.
      $\Gamma \Entails !B$ હોવાથી $\mSat{M}{!B}[w]$.
      એ જ રીતે $\mSat{M}{!C}[w]$. તેથી
      $\mSat{M}{!B \land !C}[w]$.

  \item નિષ્પત્તિનો છેલ્લો નિયમ $\Elim{\land}$ છે:
    પૂર્વધારણા
      $!B \land !C$નું નિષ્પત્તિ, જેમાં $\Gamma$ની
      ધારણાઓ અનિવૃત્ત છે, $\Gamma \Proves !B \land !C$
      બતાવે છે. આગમનની ધારણાથી $\Gamma \Entails !B \land !C$.
      $\Gamma \Entails !B$ બતાવવાનું છે. તેથી ધારો કે
      $\mSat{M}{\Gamma}[w]$. $\Gamma \Entails !B \land !C$
      હોવાથી $\mSat{M}{!B \land !C}[w]$; તેથી
      $\mSat{M}{!B}[w]$ પણ છે. એ જ રીતે જો $\Elim\land$નો
      નિષ્કર્ષ~$!C$ હોય, તો $\Gamma \Entails !C$.

  \item નિષ્પત્તિનો છેલ્લો નિયમ $\Intro{\lor}$ છે:
    ધારો કે પૂર્વધારણા $!B$
      છે અને $!B$ પર પૂરાં થતા નિષ્પત્તિની
      અનિવૃત્ત ધારણાઓ $\Gamma$ છે. ત્યારે
      $\Gamma \Proves !B$ અને આગમનની ધારણાથી
      $\Gamma \Entails !B$. $\Gamma \Entails !B \lor !C$
      બતાવવાનું છે. ધારો કે $\mSat{M}{\Gamma}[w]$.
      $\Gamma \Entails !B$ હોવાથી $\mSat{M}{!B}[w]$;
      તેથી $\mSat{M}{!B \lor !C}[w]$ પણ છે. એ જ રીતે
      પૂર્વધારણા~$!C$ હોય, તો $\Gamma \Entails !C$.

  \item નિષ્પત્તિનો છેલ્લો નિયમ~$\Elim{\lor}$ છે:
    પૂર્વધારણાઓ પર પૂરાં થતાં
      નિષ્પત્તિઓ અનુક્રમે અનિવૃત્ત ધારણાઓ~$\Gamma$
      પરથી $!B \lor !C$નાં, અનિવૃત્ત ધારણાઓ
      $\Delta_1 \cup \{!B\}$ પરથી $!D$નાં અને
      અનિવૃત્ત ધારણાઓ~$\Delta_2 \cup \{!C\}$
      પરથી $!D$નાં છે. તેથી $\Gamma \Proves !B \lor !C$,
      $\Delta_1 \cup \{!B\} \Proves !D$, અને
      $\Delta_2 \cup \{!C\} \Proves !D$.
      આગમનની ધારણાથી $\Gamma \Entails !B \lor !C$,
      $\Delta_1 \cup \{!B\} \Entails !D$, અને
      $\Delta_2 \cup \{!C\} \Entails !D$.
      $\Gamma \cup \Delta_1 \cup \Delta_2 \Entails !D$
      સાબિત કરવાનું છે.

    ધારો કે $\mSat{M}{\Gamma \cup \Delta_1 \cup \Delta_2}[w]$.
    ત્યારે $\mSat{M}{\Gamma}[w]$ છે અને
    $\Gamma \Entails !B \lor !C$ હોવાથી
    $\mSat{M}{!B \lor !C}[w]$. $\mSat{M}{}$ની
    વ્યાખ્યા મુજબ $\mSat{M}{!B}[w]$ અથવા
    $\mSat{M}{!C}[w]$. તેથી બે કિસ્સા છે:
    (a)~$\mSat{M}{!B}[w]$. ત્યારે
    $\mSat{M}{\Delta_1 \cup \{!B\}}[w]$.
    \sourcecorrection{OLMD-091}{મૂળમાં (a)માં
    $\mSat{M}{!B}$ પછીનું $[w]$ ગણિત-મર્યાદાની
    બહાર છે; તેને એ જ સત્ય-સંજ્ઞામાં મૂક્યું છે.}
    $\Delta_1 \cup \{!B\} \Entails !D$ હોવાથી
    $\mSat{M}{!D}[w]$. (b)~$\mSat{M}{!C}[w]$.
    ત્યારે $\mSat{M}{\Delta_2 \cup \{!C\}}[w]$.
    $\Delta_2 \cup \{!C\} \Entails !D$ હોવાથી
    $\mSat{M}{!D}[w]$. આથી બંને કિસ્સામાં ઇચ્છિત
    $\mSat{M}{!D}[w]$ મળે છે.
    \sourcecorrection{OLMD-092}{મૂળના બે અનુસરણ-દાવામાં
    $!B$ અને $!C$ પહેલાં $\cup$ આવે છે, પરંતુ
    એકલ સૂત્રો ગણો નથી; અગાઉની પૂર્વધારણાઓ પ્રમાણે
    તેમને $\{!B\}$ અને $\{!C\}$ એકક ગણો કર્યા છે.}

  \item નિષ્પત્તિનો છેલ્લો નિયમ $\Intro{\lif}$ છે અને
    નિષ્કર્ષ~$!B\lif !C$ છે. ત્યારે પૂર્વધારણા~$!C$
    છે અને તેના પર પૂરાં થતા નિષ્પત્તિની
    અનિવૃત્ત ધારણાઓ~$\Gamma \cup \{!B\}$ છે.
    તેથી $\Gamma \cup \{!B\} \Proves !C$ અને આગમનની
    ધારણાથી $\Gamma \cup \{!B\} \Entails !C$.
    $\Gamma \Entails !B \lif !C$ બતાવવાનું છે.

    ધારો કે $\mSat{M}{\Gamma}[w]$. દરેક એવા $w'$ માટે
    કે $Rww'$, જો $\mSat{M}{!B}[w']$ હોય, તો
    $\mSat{M}{!C}[w']$ છે તે બતાવવું છે. તેથી
    $Rww'$ અને $\mSat{M}{!B}[w']$ ધારો.
    \olref[sem][rel]{prop:true-monotonic} મુજબ
    $\mSat{M}{\Gamma}[w']$. હવે
    $\Gamma \cup \{!B\} \Entails !C$ હોવાથી
    $\mSat{M}{!C}[w']$ મળે છે, જે બતાવવાનું હતું.

  \item નિષ્પત્તિનો છેલ્લો નિયમ $\Elim{\lif}$ છે અને
    નિષ્કર્ષ~$!C$ છે. પૂર્વધારણાઓ $!B \lif !C$ અને
    $!B$ છે, જેમના નિષ્પત્તિઓ અનુક્રમે
    અનિવૃત્ત ધારણાઓ $\Gamma$ અને $\Delta$
    પરથી છે. તેથી $\Gamma \Proves !B \lif !C$ અને
    $\Delta \Proves !B$. આગમનની ધારણાથી
    $\Gamma \Entails !B \lif !C$ અને
    $\Delta \Entails !B$. $\Gamma \cup \Delta
    \Entails !C$ બતાવવાનું છે.

    ધારો કે $\mSat{M}{\Gamma \cup \Delta}[w]$.
    $\mSat{M}{\Gamma}[w]$ અને
    $\Gamma \Entails !B \lif !C$ હોવાથી
    $\mSat{M}{!B \lif !C}[w]$. વ્યાખ્યા મુજબ,
    $Rww'$ ધરાવતા દરેક $w'$ માટે જો
    $\mSat{M}{!B}[w']$ હોય, તો
    $\mSat{M}{!C}[w']$. $R$ સ્વવાચક હોવાથી
    $Rww'$ ધરાવતા $w'$માં $w$ પોતે પણ આવે છે;
    એટલે જો $\mSat{M}{!B}[w]$ હોય, તો
    $\mSat{M}{!C}[w]$. હવે $\mSat{M}{\Delta}[w]$
    અને $\Delta \Entails !B$ હોવાથી
    $\mSat{M}{!B}[w]$. તેથી ઇચ્છિત
    $\mSat{M}{!C}[w]$ મળે છે.

  \item નિષ્પત્તિનો છેલ્લો નિયમ $\FalseInt$ છે અને
    નિષ્કર્ષ~$!A$ છે. પૂર્વધારણા $\lfalse$ છે અને તેના
    નિષ્પત્તિની અનિવૃત્ત ધારણાઓ~$\Gamma$ છે.
    ત્યારે $\Gamma \Proves \lfalse$. આગમનની ધારણાથી
    $\Gamma \Entails \lfalse$. હવે $\Gamma \Entails !A$
    બતાવવાનું છે.

    પરોક્ષ રીતે આગળ વધીએ. જો $\Gamma \Entails/ !A$
    હોય, તો એવું નિદર્શ~$\mModel{M}$ અને જગત~$w$
    છે કે $\mSat{M}{\Gamma}[w]$ અને
    $\mSat/{M}{!A}[w]$. $\Gamma \Entails \lfalse$
    હોવાથી $\mSat{M}{\lfalse}[w]$. પરંતુ આ અશક્ય
    છે, કારણ કે વ્યાખ્યા મુજબ $\mSat/{M}{\lfalse}[w]$.
    તેથી $\Gamma \Entails !A$.
  \item નિષ્પત્તિનો છેલ્લો નિયમ $\Intro\lnot$ છે: સ્વાધ્યાય.
  \item નિષ્પત્તિનો છેલ્લો નિયમ $\Elim\lnot$ છે: સ્વાધ્યાય.
  \end{enumerate}
\end{proof}

\begin{prob}
  \olref[int][sc][snd]{thm:soundness}ની સાબિતી પૂરી કરો.
  $\Intro\lnot$ અને $\Elim\lnot$ના કિસ્સાઓમાં
  \olref[int][sem][rel]{defn:true-at-w}માં આપેલી
  $\mSat{M}{\lnot !A}[w]$ની વ્યાખ્યા વાપરો;
  એટલે $\lnot !A$ને $!A \lif \lfalse$ વડે
  વ્યાખ્યાયિત થયેલું માનીને કામ ન કરો.
\end{prob}

\begin{prob}
  નીચેના સૂત્રો સ્ફુરણાવાદી તર્કમાં નિષ્પન્ન કરી શકાય એવું
  નથી તે બતાવો:
  \begin{enumerate}
    \item $(!A \lif !B) \lor (!B \lif !A)$
    \item $(\lnot\lnot !A \lif !A) \lif (!A \lor \lnot !A)$
    \item $(!A \lif !B \lor !C) \lif \bigl((!A \lif !B)\lor(!A \lif !C)\bigr)$
  \end{enumerate}
\end{prob}
% END OLP-0505
\endgroup

\begingroup
\makeatletter\def\ol@sectioncs{section}\makeatother

% BEGIN OLP-0506
\olfileid{int}{sc}{lin}

\olsection{લિન્ડનબાઉમનું સહાયક પ્રમેય}
\phantomsection\label{guunit:OLP-0506}


સ્ફુરણાવાદી તર્ક માટેનો પૂર્ણતા પ્રમેય $\Gamma \Proves/ !A$
માનીને $\mSat{M}{\Gamma}$ અને~$\mSat/{M}{!A}$ ધરાવતું
નિદર્શ રચીને સાબિત થાય છે.

પ્રશિષ્ટ તર્કમાં નિષ્પન્નક્ષમતાના સંબંધને સુસંગતતાના
ખ્યાલમાં ઉતારી શકાય છે: સૂત્ર~$!A$ એ
સૂત્રોના ગણમાંથી ત્યારે અને માત્ર ત્યારે જ
નિષ્પન્ન કરી શકાય એવું હોય છે જ્યારે તે ગણ અને~$!A$નો નિષેધ
મળીને અસુસંગત થાય. સ્ફુરણાવાદી તર્કમાં આવું શક્ય નથી.
ત્યાં જો $\lnot!A$ અસુસંગત હોય, તો માત્ર
$\Proves \lnot\lnot !A$ મળે છે. સામાન્ય રીતે
$\lnot\lnot!A \lif !A$ સ્ફુરણાવાદી રીતે માન્ય નથી,
તેથી $\Proves !A$ તારવી શકાતું નથી.

આથી નિદર્શ~$\mModel{M}$ રચતી વખતે સૂત્ર~$!A$નું
નિષ્પન્નક્ષમતા ન હોવું ધ્યાનમાં રાખવું પડે છે. તેથી નિદર્શ
$\mModel{M}$ રચવા માટે પૂર્ણ ગણ $\Gamma^* \supseteq \Gamma$
વાપરી શકાતો નથી, કારણ કે દરેક પૂર્ણ ગણ~$\Gamma^*$માં
$\Gamma^* \Proves !A \lor \lnot !A$ હોય છે.

પૂર્ણ ગણ~$\Gamma^*$ને બદલે આપણે સૂત્રોના અભાજ્ય ગણનો
ખ્યાલ વાપરીશું:

\begin{defn}\ollabel{defn:prime}
  સૂત્રોનો ગણ~$\Gamma$ ત્યારે અને માત્ર ત્યારે જ
  \emph{અભાજ્ય} કહેવાય જ્યારે
  \begin{enumerate}
  \item\ollabel{defn:prime1} $\Gamma$ સુસંગત હોય, એટલે
    $\Gamma \Proves/ \lfalse$;
  \item\ollabel{defn:prime2} જો $\Gamma \Proves !A$ હોય,
    તો $!A \in \Gamma$; અને
  \item\ollabel{defn:prime3} જો $!A \lor !B \in \Gamma$
    હોય, તો $!A \in \Gamma$ અથવા $!B \in \Gamma$.
  \end{enumerate}
\end{defn}

\begin{lem}[લિન્ડનબાઉમનું સહાયક પ્રમેય]
  \ollabel{lem:lindenbaum} જો $\Gamma \Proves/ !A$ હોય,
  તો એવો $\Gamma^* \supseteq \Gamma$ છે કે $\Gamma^*$
  અભાજ્ય છે અને $\Gamma^* \Proves/ !A$.
\end{lem}

\begin{proof}
  $!B_1 \lor !C_1$, $!B_2 \lor !C_2$, \dots, એ
  $!B \lor !C$ સ્વરૂપના બધા સૂત્રોની પરિગણના હોય.
  સૂત્રોના ગણોની વધતી શ્રેણી~$\Gamma_n$ વ્યાખ્યાયિત
  કરીએ છીએ, જેમાં દરેક $\Gamma_{n+1}$ એ $\Gamma_n$માં
  એક નવું સૂત્ર ઉમેરીને બને છે. બધા~$\Gamma_n$નો
  યોગગણ $\Gamma^*$ હશે. નવા સૂત્રો એવી રીતે પસંદ
  કરીએ છીએ કે $\Gamma^*$ અભાજ્ય રહે અને છતાં
  $\Gamma^* \Proves/ !A$ હોય. એટલે દરેક તબક્કે
  નીચેની શરતો સંતોષતું પહેલું વિયોજન~$!B_i \lor !C_i$
  શોધવું પડશે:
  \begin{enumerate}
  \item\ollabel{gamma-1} $\Gamma_n \Proves !B_i \lor !C_i$
  \item\ollabel{gamma-2} $!B_i \notin \Gamma_n$ અને
    $!C_i \notin \Gamma_n$
  \end{enumerate}
  જો $\Gamma_n \cup \{!B_i\} \Proves/ !A$ હોય, તો
  $\Gamma_n$માં $!B_i$ ઉમેરીએ; નહિ તો $!C_i$ ઉમેરીએ.
  આ રીત કામ કરે છે તે બતાવવું પડશે. હાલ માટે
  \olref{gamma-1} અને~\olref{gamma-2} સંતોષતા સૌથી
  નાના $i$ને $i(n)$ કહીએ.

  $\Gamma_0 = \Gamma$ મૂકો અને
  \[
  \Gamma_{n+1} = \begin{cases}
    \Gamma_n \cup \{!B_{i(n)}\} &
    \text{જો $\Gamma_n \cup \{!B_{i(n)}\} \Proves/ !A$ હોય} \\
    \Gamma_n \cup \{!C_{i(n)}\} & \text{અન્યથા}
    \end{cases}
  \]
  જો $i(n)$ વ્યાખ્યાયિત ન હોય, એટલે $\Gamma_n \Proves !B \lor !C$
  થાય ત્યારે હંમેશાં $!B \in \Gamma_n$ અથવા $!C \in \Gamma_n$
  હોય, તો $\Gamma_{n+1} = \Gamma_n$ મૂકો.
  હવે $\Gamma^* = \bigcup_{n=0}^\infty \Gamma_n$ મૂકો.

  પહેલાં દરેક $n$ માટે $\Gamma_n \Proves/ !A$ બતાવીએ.
  $n$ પર આગમન કરીએ. $n = 0$ માટે સહાયક પ્રમેયની ધારણા,
  એટલે $\Gamma \Proves/ !A$, પરથી દાવો મળે છે.
  આગમનના પગલામાં કોઈ પણ $n\ge 0$ માટે
  $\Gamma_n \Proves/ !A$ પરથી
  $\Gamma_{n+1} \Proves/ !A$ બતાવવું છે. $i(n)$
  વ્યાખ્યાયિત ન હોય, તો $\Gamma_{n+1} = \Gamma_n$
  હોવાથી કશું સાબિત કરવાનું નથી. તેથી $i(n)$
  વ્યાખ્યાયિત છે એમ ધારો. સરળતા માટે $i = i(n)$ મૂકો.
  \sourcecorrection{OLMD-093}{મૂળમાં આધાર $n=0$ પછી
  આગમનનું પગલું $n>0$થી શરૂ થાય છે, જેથી
  $\Gamma_0$થી $\Gamma_1$નો કિસ્સો છૂટી જાય છે.
  પગલું બધા $n\ge 0$ માટે જણાવ્યું છે.}

  દાવાનું પ્રતિપક્ષી વિધાન સાબિત કરીએ. ધારો કે
  $\Gamma_{n+1} \Proves !A$. રચના મુજબ જો
  $\Gamma_n \cup \{!B_i\} \Proves/ !A$ હોય, તો
  $\Gamma_{n+1} = \Gamma_n \cup \{!B_i\}$; નહિ તો
  $\Gamma_{n+1} = \Gamma_n \cup \{!C_i\}$.
  પહેલો કિસ્સો બની શકે નહીં, કારણ કે ત્યારે
  $\Gamma_{n+1} \Proves/ !A$ થાય. તેથી
  $\Gamma_n \cup \{!B_i\} \Proves !A$ અને
  $\Gamma_{n+1} = \Gamma_n \cup \{!C_i\}$.
  $i(n)$ની વ્યાખ્યાથી $\Gamma _n \Proves !B_i \lor !C_i$.
  વળી $\Gamma_n \cup \{!B_i\} \Proves !A$ અને
  $\Gamma_{n+1} = \Gamma_n \cup \{!C_i\} \Proves !A$.
  આથી ઇચ્છિત $\Gamma_n \Proves !A$ મળે છે.

  જો $\Gamma^* \Proves !A$ હોય, તો એવો સાન્ત ઉપગણ
  $\Gamma' \subseteq \Gamma^*$ હોય
  કે $\Gamma' \Proves !A$. તેનું દરેક
  $!D \in \Gamma'$ કોઈક~$i$ માટે $\Gamma_i$માં હશે.
  આ સૂચકોમાં સૌથી મોટો $n$ લો. $i \le n$ હોય ત્યારે
  $\Gamma_i \subseteq \Gamma_n$ હોવાથી
  $\Gamma' \subseteq \Gamma_n$. તેથી
  $\Gamma_n \Proves !A$, જે ઉપર સાબિત કરેલા
  $\Gamma_n \Proves/ !A$થી વિરુદ્ધ છે.

  છેલ્લે $\Gamma^*$ અભાજ્ય છે, એટલે કે
  \olref{defn:prime}ની \olref{defn:prime1},
  \olref{defn:prime2} અને~\olref{defn:prime3}
  શરતો સંતોષે છે તે બતાવીએ.

  સૌપ્રથમ $\Gamma^* \Proves/ !A$ છે, તેથી
  $\Gamma^*$ સુસંગત છે અને \olref{defn:prime1} મળે છે.

  હવે જો $\Gamma^* \Proves !B \lor !C$ હોય, તો
  $!B \in \Gamma^*$ અથવા $!C \in \Gamma^*$ છે તે
  બતાવીએ. આનાથી \olref{defn:prime3} મળે છે, કારણ કે
  $!B \lor !C \in \Gamma^*$ હોય, તો
  $\Gamma^* \Proves !B \lor !C$ પણ હોય. તેથી
  $\Gamma^* \Proves !B \lor !C$, પરંતુ
  $!B \notin \Gamma^*$ અને $!C \notin \Gamma^*$,
  એમ ધારો. $\Gamma^* \Proves !B \lor !C$
  હોવાથી કોઈક~$n$ માટે $\Gamma_n \Proves !B \lor !C$.
  બધા વિયોજનોની પરિગણનામાં $!B \lor !C$ કોઈક
  $!B_j \lor !C_j$ તરીકે આવે છે. $!B_j \lor !C_j$
  એ $i(n)$ની વ્યાખ્યાની શરતો સંતોષે છે: એટલે
  $\Gamma_n \Proves !B_j \lor !C_j$, જ્યારે
  $!B_j \notin \Gamma_n$ અને $!C_j \notin \Gamma_n$.
  દરેક તબક્કે પસંદ કરેલું વિયોજન~$!B_i \lor !C_i$
  પછીની યાદીમાં ફરી શરતો સંતોષતું નથી, કારણ કે
  તેમાં $!B_i$ અથવા $!C_i$ ઉમેરાય છે. આ સૂચક
  પહેલાં સાન્ત સૂચકો જ છે; તેથી આ વિયોજન શરતો સંતોષતું
  રહે ત્યાં સુધી તેનાથી નાના બાકી સૂચકો એક પછી એક
  દૂર થશે અને કોઈક તબક્કે~$m$ માટે $j = i(m)$ મળશે.
  \sourcecorrection{OLMD-094}{મૂળ કહે છે કે દરેક તબક્કે
  શરતો સંતોષતા વિયોજનોની કુલ સંખ્યા ઓછામાં ઓછી એકથી
  ઘટે છે; નવા વિયોજનો નિગમ્ય બની શકે હોવાથી આ દાવો
  જરૂરી નથી. નક્કી કરેલા $j$થી નાના સાન્ત સૂચકો
  ફરી પસંદ ન થાય એ તર્ક વાપર્યો છે.}
  પરંતુ પછી $!B \in \Gamma_{m+1}$ અથવા
  $!C \in \Gamma_{m+1}$, જે $!B \notin \Gamma^*$
  અને $!C \notin \Gamma^*$ની ધારણાથી વિરુદ્ધ છે.

  હવે ધારો કે $\Gamma^* \Proves !B$. તો
  $\Gamma^* \Proves !B \lor !B$. હમણાં સાબિત કર્યું
  છે કે $\Gamma^* \Proves !B \lor !B$ હોય, તો
  $!B \in \Gamma^*$. તેથી $\Gamma^*$ એ
  \olref{defn:prime}ની \olref{defn:prime2} સંતોષે છે.
\end{proof}

\begin{prob}
જો $\Gamma \Proves/ \bot$ હોય, તો પ્રશિષ્ટ તર્કમાં
$\Gamma$ સુસંગત છે, એટલે~$\Gamma$નાં બધાં સૂત્રોને
સાચાં બનાવતું મૂલ્યાંકન છે, તે બતાવો.
\end{prob}
% END OLP-0506
\endgroup

\begingroup
\makeatletter\def\ol@sectioncs{section}\makeatother

% BEGIN OLP-0507
\olfileid{int}{sc}{mod}

\olsection{કૅનોનિકલ નિદર્શ}
\phantomsection\label{guunit:OLP-0507}


આપણા નિદર્શનાં જગતો પ્રાકૃતિક સંખ્યાઓની સાન્ત
શ્રેણીઓ~$\sigma$ હશે, એટલે $\sigma \in \Nat^*$.
નોંધો કે $\Nat^*$ને આગમનાત્મક રીતે આ પ્રમાણે વ્યાખ્યાયિત
કરીએ છીએ:
\begin{enumerate}
\item $\emptyseq \in \Nat^*$.
\item જો $\sigma \in \Nat^*$ અને $n \in \Nat$ હોય,
  તો $\sigma.n \in \Nat^*$ (અહીં $\sigma.n$ એ
  $\sigma \concat \tuple{n}$ છે અને $\sigma \concat \sigma'$
  એ $\sigma$ તથા $\sigma'$નું જોડાણ છે).
\item $\Nat^*$માં બીજું કશું નથી.
\end{enumerate}
તેથી આગમનાત્મક વ્યાખ્યાઓ આપવા $\Nat^*$ વાપરી શકીએ છીએ.

$\tuple{!B_1, !C_1}$, $\tuple{!B_2, !C_2}$, \dots,
એ સૂત્રોની બધી જોડીઓની પરિગણના હોય.
સૂત્રોનો ગણ~$\Delta$ આપેલો હોય, તો
$\Delta(\sigma)$ને નીચે મુજબ આગમનાત્મક રીતે
વ્યાખ્યાયિત કરીએ છીએ:
\begin{enumerate}
\item $\Delta(\emptyseq) = \Delta$
\item $\Delta(\sigma.n) = {}$
  \[
  \begin{cases}
    (\Delta(\sigma) \cup \{!B_n\})^* &
    \text{જો $\Delta(\sigma) \cup \{!B_n\} \Proves/ !C_n$ હોય} \\
    \Delta(\sigma) & \text{અન્યથા}
  \end{cases}
  \]
\end{enumerate}
અહીં $(\Delta(\sigma) \cup \{!B_n\})^*$નો અર્થ
સૂત્રોનો એવો અભાજ્ય ગણ છે, જે
$\Delta(\sigma) \cup \{!B_n\}$ અને સૂત્ર~$!C_n$
પર \olref[lin]{lem:lindenbaum} લાગુ કરતાં મળે છે.
આ વ્યાખ્યા મુજબ જો $\Delta(\sigma) \cup \{!B_n\}
\Proves/ !C_n$ હોય, તો $\Delta(\sigma.n)
\Proves !B_n$ અને $\Delta(\sigma.n) \Proves/ !C_n$.
એ પણ નોંધો કે કોઈ પણ~$n$ માટે
$\Delta(\sigma) \subseteq \Delta(\sigma.n)$.
જો $\Delta$ અભાજ્ય હોય, તો દરેક~$\sigma$ માટે
$\Delta(\sigma)$ અભાજ્ય હોય છે.

\begin{defn}\ollabel{defn:canonical-model}
  ધારો કે $\Delta$ અભાજ્ય છે. ત્યારે $\Delta$ માટેનું
  \emph{કૅનોનિકલ નિદર્શ} $\mModel{M(\Delta)}$ આ પ્રમાણે
  વ્યાખ્યાયિત થાય છે:
  \begin{enumerate}
  \item $W = \Nat^*$, એટલે પ્રાકૃતિક સંખ્યાઓની સાન્ત
    શ્રેણીઓનો ગણ.
  \item $R$ એવો આંશિક ક્રમ છે કે $R\sigma\sigma'$
    ત્યારે અને માત્ર ત્યારે જ જ્યારે $\sigma$ એ~$\sigma'$નો
    પ્રારંભિક ખંડ હોય (એટલે કોઈ શ્રેણી~$\sigma''$ માટે
    $\sigma' = \sigma \concat \sigma''$ હોય).
  \item $V(p) = \Setabs{\sigma}{p \in \Delta(\sigma)}$.
  \end{enumerate}
\end{defn}

$R$ ખરેખર આંશિક ક્રમ છે તે સહેલાઈથી ચકાસી શકાય છે.
$V$ માટેની એકદિશવર્ધિતાની શરત પણ સંતોષાય છે.
$\Delta(\sigma) \subseteq \Delta(\sigma.n)$ હોવાથી
$R\sigma\sigma'$ હોય ત્યારે $\sigma$ પછી જોડાતા ભાગની
લંબાઈ પર આગમન કરીને
$\Delta(\sigma) \subseteq \Delta(\sigma')$ મળે છે.
\sourcecorrection{OLMD-095}{મૂળ આ છેલ્લું આગમન
$\sigma$ પર કહે છે; $\sigma$ સ્થિર રાખીને તેના પછી
જોડાતી સાન્ત શ્રેણીની લંબાઈ પર આગમન જોઈએ.}
% END OLP-0507
\endgroup

\begingroup
\makeatletter\def\ol@sectioncs{section}\makeatother

% BEGIN OLP-0508
\olfileid{int}{sc}{tru}

\olsection{સત્યતા સહાયક પ્રમેય}
\phantomsection\label{guunit:OLP-0508}


નીચેનું સહાયક પ્રમેય કૅનોનિકલ નિદર્શમાં સત્ય હોવાને
કયા સૂત્રો અભાજ્ય ગણ~$\Delta$ના ઘટકો
છે તેની સાથે જોડે છે.

\begin{lem}\ollabel{lem:truth}
  જો $\Delta$ અભાજ્ય હોય, તો $\mSat{M(\Delta)}{!A}[\sigma]$
  ત્યારે અને માત્ર ત્યારે જ જ્યારે
  $\Delta(\sigma) \Proves !A$.
\end{lem}

\begin{proof}
  $!A$ પર આગમન કરીએ.
  \begin{enumerate}
    \item \indcase{!A}{\lfalse}{$\Delta(\sigma)$ અભાજ્ય છે,
      તેથી સુસંગત છે અને $\Delta(\sigma) \Proves/ \indfrm$.
      વ્યાખ્યા મુજબ $\mSat/{M(\Delta)}{\indfrm}[\sigma]$.}
      \item \indcase{!A}{p}{$\mSat{{}}{}$ની વ્યાખ્યા મુજબ
        $\mSat{M(\Delta)}{\indfrm}[\sigma]$ ત્યારે અને માત્ર
        ત્યારે જ જ્યારે $\sigma \in V(p)$, એટલે
        $\Delta(\sigma) \Proves \indfrm$.}
      \item \indcase!{!A}{\lnot !B}{}
      \item \indcase{!A}{!B \land
        !C}{$\mSat{M(\Delta)}{\indfrm}[\sigma]$ ત્યારે અને
        માત્ર ત્યારે જ જ્યારે
        $\mSat{M(\Delta)}{!B}[\sigma]$ અને
        $\mSat{M(\Delta)}{!C}[\sigma]$. આગમનની ધારણાથી
        $\mSat{M(\Delta)}{!B}[\sigma]$ ત્યારે અને માત્ર
        ત્યારે જ જ્યારે $\Delta(\sigma) \Proves
        !B$; $!C$ માટે પણ એ જ રીતે. હવે
        $\Delta(\sigma) \Proves !B$ અને
        $\Delta(\sigma) \Proves !C$ ત્યારે અને માત્ર ત્યારે
        જ જ્યારે $\Delta(\sigma)
        \Proves \indfrm$.}
      \item \indcase{!A}{!B \lor
        !C}{$\mSat{M(\Delta)}{\indfrm}[\sigma]$ ત્યારે અને
        માત્ર ત્યારે જ જ્યારે
        $\mSat{M(\Delta)}{!B}[\sigma]$ અથવા
        $\mSat{M(\Delta)}{!C}[\sigma]$. આગમનની ધારણાથી
        આ ત્યારે અને માત્ર ત્યારે જ જ્યારે
        $\Delta(\sigma) \Proves !B$ અથવા
        $\Delta(\sigma) \Proves !C$. હવે આ શરત
        $\Delta(\sigma) \Proves \indfrm$ સાથે સમકક્ષ
        છે તે બતાવવું છે. ડાબેથી જમણી દિશા સ્પષ્ટ છે.
        જમણેથી ડાબી દિશા $\Delta(\sigma)$ અભાજ્ય
        હોવાથી મળે છે.}
      \item \indcase{!A}{!B \lif !C}{પહેલાં ડાબેથી જમણી
        દિશાનું પ્રતિપક્ષી વિધાન. ધારો કે
        $\Delta(\sigma) \Proves/ !B
        \lif !C$. તો $\Delta(\sigma) \cup \{!B\} \Proves/
        !C$ પણ છે. કોઈક~$n$ માટે $\tuple{!B, !C}$ એ
        $\tuple{!B_n, !C_n}$ હોવાથી
        $\Delta(\sigma.n) = (\Delta(\sigma) \cup
        \{!B\})^*$, અને $\Delta(\sigma.n) \Proves !B$
        પરંતુ $\Delta(\sigma.n) \Proves/ !C$. આગમનની
        ધારણાથી $\mSat{M(\Delta)}{!B}[\sigma.n]$
        અને $\mSat/{M(\Delta)}{!C}[\sigma.n]$.
        $R\sigma(\sigma.n)$ હોવાથી આનો અર્થ
        $\mSat/{M(\Delta)}{\indfrm}[\sigma]$ થાય છે.

        હવે ધારો કે $\Delta(\sigma) \Proves !B \lif !C$
        અને $R\sigma\sigma'$. $\Delta(\sigma) \subseteq
        \Delta(\sigma')$ હોવાથી, જો
        $\Delta(\sigma') \Proves !B$ હોય, તો
        $\Delta(\sigma') \Proves !C$. બીજા શબ્દોમાં,
        $R\sigma\sigma'$ ધરાવતા દરેક $\sigma'$ માટે
        $\Delta(\sigma') \Proves/ !B$ અથવા
        $\Delta(\sigma') \Proves !C$. આગમનની ધારણાથી
        $R\sigma\sigma'$ હોય ત્યારે
        $\mSat/{M(\Delta)}{!B}[\sigma']$ અથવા
        $\mSat{M(\Delta)}{!C}[\sigma']$;
        એટલે $\mSat{M(\Delta)}{\indfrm}[\sigma]$.}
  \end{enumerate}
\end{proof}
% END OLP-0508
\endgroup

\begingroup
\makeatletter\def\ol@sectioncs{section}\makeatother

% BEGIN OLP-0509
\olfileid{int}{sc}{cpl}

\olsection{પૂર્ણતા પ્રમેય}
\phantomsection\label{guunit:OLP-0509}


\begin{thm}\ollabel{thm:completeness}
  જો $\Gamma \Entails !A$ હોય, તો $\Gamma \Proves !A$.
\end{thm}

\begin{proof}
  પ્રતિપક્ષી વિધાન સાબિત કરીએ: ધારો કે $\Gamma \Proves/ !A$.
  ત્યારે \olref[lin]{lem:lindenbaum} મુજબ એવો અભાજ્ય ગણ
  $\Gamma^* \supseteq \Gamma$ છે કે $\Gamma^* \Proves/ !A$.
  \olref[mod]{defn:canonical-model}માં વ્યાખ્યાયિત
  $\Gamma^*$ માટેનું કૅનોનિકલ નિદર્શ~$\mModel{M(\Gamma^*)}$
  લો. દરેક $!B \in \Gamma$ માટે $\Gamma^* \Proves !B$.
  નોંધો કે $\Gamma^*(\emptyseq) = \Gamma^*$.
  સત્યતા સહાયક પ્રમેય (\olref[tru]{lem:truth}) મુજબ
  બધા $!B \in \Gamma$ માટે
  $\mSat{M(\Gamma^*)}{!B}[\emptyseq]$ અને
  $\mSat/{M(\Gamma^*)}{!A}[\emptyseq]$ મળે છે.
  આથી $\Gamma \Entails/ !A$.
\end{proof}

\begin{prob}
  જો $!A$માં ફક્ત પ્રતિજ્ઞાપન ચલો, $\lor$
  અને~$\land$ હોય, તો $\Entails/ !A$ છે તે બતાવો.
  આ પરથી સ્ફુરણાવાદી તર્કમાં $\lor$ અને~$\land$
  વડે $\lif$ વ્યાખ્યાયિત થઈ શકતું નથી તે તારવો.
\end{prob}

\begin{prob}
  પૂર્ણતા પ્રમેય વાપરીને સાબિત કરો કે જો
  $\Proves !A \lor !B$ હોય, તો $\Proves !A$ અથવા
  $\Proves !B$. (સંકેત: $\mSat/{M_1}{!A}$ અને
  $\mSat/{M_2}{!B}$ ધારો અને એવું નવું નિદર્શ
  $\mModel{M}$ રચો કે $\mSat/{M}{!A \lor !B}$.)
\end{prob}

\begin{prob}
  જો $\mModel{M}$ રેખીય ક્રમ વાપરતું સંબંધાત્મક
  નિદર્શ હોય, તો
  $\mSat{M}{(!A \lif !B)\lor(!B \lif !A)}$
  છે તે બતાવો.
\end{prob}
% END OLP-0509
\endgroup

\begingroup
\makeatletter\def\ol@sectioncs{section}\makeatother

% BEGIN OLP-0510
\olfileid{int}{sc}{dec}

\olsection{નિર્ણેયતા}
\phantomsection\label{guunit:OLP-0510}

નોંધો કે પૂર્ણતા પ્રમેયની સાબિતી દરેક
$\Gamma \Proves/ !A$ માટે અનંત સંખ્યામાં જગતો ધરાવતું
એવું નિદર્શ આપે છે જે $\Gamma \Entails/ !A$ની સાક્ષી
પૂરે છે. નીચેનું વિધાન બતાવે છે કે $\Entails !A$
સાબિત કરવા બધા સાન્ત નિદર્શો માટે
$\mSat{M}{!A}$ સાબિત કરવું પૂરતું છે (એટલે, સાન્ત
જગતગણ ધરાવતા નિદર્શો માટે).

\begin{thm}\ollabel{thm:decidability}
  જો $\Entails/ !A$ હોય, તો એવું સાન્ત નિદર્શ છે કે
  $\mSat/{M'}{!A}$.
\end{thm}

\begin{proof}
  ધારો કે $\mModel{M}=\tuple{W, R, V}$ એવું છે કે
  $\mSat/{M}{!A}$ અને $P$ એ~$!A$માં આવતા
  પ્રતિજ્ઞાપન ચલોનો ગણ છે.
  $W' = \Setabs{[w]}{w \in W}$ લઈને, જ્યાં
  $[w] = \Setabs{p \in P}{w \in V(p)}$, $R'$ને
  ઉપગણ-સંબંધ અને
  $V'(p)=\Setabs{[w]}{p \in [w]}$ રાખીને
  $\mModel{M'}=\tuple{W', R', V'}$ વ્યાખ્યાયિત કરો.
  $W'$ સાન્ત ગણ છે અને $\mModel{M'}$ એ
  સંબંધી નિદર્શ છે તે સ્પષ્ટ છે.

  $!A$ પર આગમન કરીને બતાવી શકાય છે કે
  \[
    \mSat{M}{!A}[w] \text{ ત્યારે અને માત્ર ત્યારે જ જ્યારે } \mSat{M'}{!A}[{[w]}]
  \]
  આ $P$માંથી જ પ્રતિજ્ઞાપન ચલો ધરાવતા
  બધા સૂત્રો~$!A$ માટે સાચું છે. તેની સાબિતી
  વાચક માટે સ્વાધ્યાય તરીકે છોડી છે.
  \sourcecorrection{OLMD-096}{મૂળમાં આપેલી આ રચના
  નિષેધ માટે સત્યતા જાળવતી નથી. દાખલા તરીકે,
  $W=\{u,v\}$, $R$ ઓળખસંબંધ અને $V(p)=\{v\}$
  હોય, તો $u$ ખાતે $\lnot p$ સાચું છે. પરંતુ નવા
  નિદર્શમાં $[u]=\emptyset \subseteq \{p\}=[v]$,
  એટલે $[u]$થી $p$ સાચું હોય એવા જગત સુધી પહોંચી
  શકાય છે અને ત્યાં $\lnot p$ ખોટું છે. સાન્ત-નિદર્શ
  પ્રમેયની બીજી સાબિતી જરૂરી છે; આ દાવાની ખામી
  નોંધ્યા વિના તેને ચકાસેલી સાબિતી માનવી નહીં.}
\end{proof}

\begin{prob}
\olref[int][sc][dec]{thm:decidability}ની સાબિતી પૂરી
કરો: ફક્ત~$P$માંથી વિધાનાત્મક ચલો ધરાવતા બધા
સૂત્રો~$!A$ માટે $\mSat{M,w}{!A}$ ત્યારે
અને માત્ર ત્યારે જ જ્યારે $\mSat{M',[w]}{!A}$
તે બતાવો.
\end{prob}

\olref{thm:decidability} પરથી $\Entails !A$ છે કે નહીં
તે નક્કી કરવાની કલનવિધિ મળે છે.
% END OLP-0510
\endgroup


\OLEndChapterHook
% END OLP-0503
\endgroup


\begingroup
\makeatletter\def\ol@sectioncs{section}\makeatother

% BEGIN OLP-0511
\olchapter{int}{tab}{સ્ફુરણાવાદી ટેબ્લો}
\olchapterid{int}{tab}
\phantomsection\label{guunit:OLP-0511}


\begin{editorial}
  સ્ફુરણાવાદી તર્ક માટે ઉપસર્ગવાળા ટેબ્લો પરનું આ પ્રારૂપ
  પ્રકરણ છે. તેમાં વધુ ઉદાહરણો, પૂર્ણતાની સાબિતીઓ અને
  બંધ ટેબ્લોની નિષ્ફળ શોધમાંથી પ્રતિનિદર્શનો કેવી રીતે
  મેળવી શકાય તેની ચર્ચા ઉમેરવાની જરૂર છે.
\end{editorial}

\begingroup
\makeatletter\def\ol@sectioncs{section}\makeatother

% BEGIN OLP-0512
\olfileid{int}{tab}{int}

\olsection{પરિચય}
\phantomsection\label{guunit:OLP-0512}


ટેબ્લો એ ચિહ્નિત સૂત્રોનાં ચોક્કસ પ્રકારનાં (નીચે તરફ શાખા પાડતાં)
વૃક્ષો છે. ચિહ્નિત સૂત્ર એટલે સત્યમૂલ્ય-ચિહ્ન
($\True$ અથવા $\False$) અને વાક્યની જોડી:
\[
\sFmla{\True}{!A} \text{ અથવા } \sFmla{\False}{!A}.
\]
ટેબ્લો કેટલીક \emph{ધારણાઓ}થી શરૂ થાય છે.
પછીનું દરેક ચિહ્નિત સૂત્ર અનુમાનનો કોઈ નિયમ
લાગુ કરીને મળે છે. કેટલાક નિયમો વૃક્ષની શાખાના છેડે
એક અથવા વધુ ચિહ્નિત સૂત્રો ઉમેરે છે; બીજા
નિયમો બે નવા છેડા ઉમેરીને બે શાખાઓ બનાવે છે.
કોઈ નિયમથી મળતા ચિહ્નિત સૂત્રોમાંનું
સૂત્ર જે ચિહ્નિત સૂત્ર પર નિયમ લાગુ થયો
તેના સૂત્ર કરતાં ઓછું સંકુલ હોય છે. કોઈ શાખામાં
$\sFmla{\True}{!A}$ અને $\sFmla{\False}{!A}$
બંને હોય, તો તેને \emph{બંધ} કહીએ છીએ.
ટેબ્લોની દરેક શાખા બંધ હોય, તો આખું
ટેબ્લો બંધ છે. બંધ ટેબ્લો એવું
નિષ્પત્તિ છે જે બતાવે છે કે ટેબ્લોની
શરૂઆતમાં લેવાયેલા
ચિહ્નિત સૂત્રોનો ગણ અસંતોષ્ય છે. આથી
$\Proves$ સંબંધ વ્યાખ્યાયિત કરી શકાય:
$\Gamma \Proves !A$ ત્યારે અને માત્ર ત્યારે જ જ્યારે
એવો કોઈ સાન્ત ગણ~$\Gamma_0 = \{!B_1,
\dots, !B_n\} \subseteq \Gamma$ હોય કે નીચેની
ધારણાઓ માટે બંધ ટેબ્લો મળે:
\[
\{\sFmla{\False}{!A}, \sFmla{\True}{!B_1}, \dots, \sFmla{\True}{!B_n}\}.
\]

સ્ફુરણાવાદી તર્ક માટે ચિહ્નિત સૂત્રની સંકલ્પના
વિસ્તૃત કરવી અને સંયોજકોના નિયમો બદલવા પડે છે.
ચિહ્ન ($\True$ અથવા $\False$) ઉપરાંત,
સ્ફુરણાવાદી ટેબ્લોનાં સૂત્રોને
\emph{ઉપસર્ગો}~$\sigma$ પણ હોય છે.
\sourcecorrection{OLMD-097}{મૂળના આ વાક્યમાં
``મોડલ ટેબ્લો'' છે, પરંતુ અહીંનું પ્રકરણ અને તેના
નિયમો સ્ફુરણાવાદી ટેબ્લો વિશે છે; વિષયનું નામ
સુસંગત કર્યું છે.}
ઉપસર્ગો ધન પૂર્ણાંકોની ખાલી ન હોય તેવી શ્રેણીઓ છે,
એટલે $\sigma \in (\PosInt)^* \setminus \{\emptyseq\}$.
આવા ઉપસર્ગો લખતી વખતે આપણે આસપાસના
$\tuple{\ }$ છોડીએ છીએ અને જુદા ઘટકો
વચ્ચે $,$ને બદલે~$.$ મૂકીએ છીએ. જો $\sigma$
ઉપસર્ગ હોય, તો $\sigma.n$ એ $\sigma \concat \tuple{n}$
છે; દાખલા તરીકે, $\sigma = 1.2.1$ હોય, તો
$\sigma.3$ એ $1.2.1.3$ છે. આમ, દાખલા તરીકે,
\[
\sFmla{\True}{!A \lif (!B \lif !C)}[1.2]
\]
એ \emph{ઉપસર્ગવાળું ચિહ્નિત સૂત્ર} છે
(અથવા ટૂંકમાં \emph{ઉપસર્ગવાળું સૂત્ર}).

સહજ રીતે, ઉપસર્ગ નિદર્શમાં એવું જગત સૂચવે છે જે
ટેબ્લોની કોઈ શાખા પરનાં સૂત્રોને
સંતોષી શકે; અને $\sigma$ કોઈ જગતનું નામ હોય, તો
$\sigma.n$ એ~$\sigma$ નામના જગતમાંથી સુલભ જગતનું
નામ આપે છે.

સ્ફુરણાવાદી નિદર્શોમાં સુલભતા સંબંધ સ્વવાચક અને
પરંપરિત હોય છે. ઉપસર્ગોની દૃષ્ટિએ તેનો અર્થ એ કે
$\sigma$ પોતે~$\sigma$માંથી સુલભ છે, અને
$\sigma$ને વિસ્તારીને મળતો કોઈ પણ ઉપસર્ગ,
એટલે $\sigma.n_1.\cdots.n_k$ સ્વરૂપનો ઉપસર્ગ,
પણ સુલભ છે. $\sigma$ પોતે અને તેનો કોઈ પણ
વિસ્તાર દર્શાવવા $\sigma.*$ સંજ્ઞા લઈએ.
બીજા શબ્દોમાં, $\sigma.*$માં બરાબર એ બધા
ઉપસર્ગો છે જે~$\sigma$માંથી સુલભ છે.
% END OLP-0512
\endgroup

\begingroup
\makeatletter\def\ol@sectioncs{section}\makeatother

% BEGIN OLP-0513
\olfileid{int}{tab}{rul}

\olsection{સ્ફુરણાવાદી તર્ક માટેના નિયમો}
\phantomsection\label{guunit:OLP-0513}


$\land$ અને $\lor$ માટેના નિયમો સામાન્ય વિધાનસંબંધી
ચિહ્નિત ટેબ્લો જેવા જ છે; ફક્ત ઉપસર્ગો ઉમેરાય છે.
દરેક કિસ્સામાં ચિહ્નિત સૂત્ર
$\sFmla{S}{!A}[\sigma]$ પર લાગુ થયેલો નિયમ નવાં
સૂત્રો આપે છે, જેમને પણ~$\sigma$ ઉપસર્ગ હોય છે.
આ સહજ છે: દાખલા તરીકે, $!A \land !B$ એ~$\sigma$
નામના જગતમાં સત્ય હોય, તો $!A$ અને $!B$ પણ
~$\sigma$માં સત્ય છે. આ નિયમના નિષ્કર્ષ પર અન્ય
જગતનો ઉપસર્ગ મૂકાતો નથી.
\sourcecorrection{OLMD-098}{મૂળનું કૌંસ કહે છે કે
$!A$ અને $!B$ બીજા કોઈ જગતમાં સત્ય નથી;
સ્ફુરણાવાદી એકદિશવર્ધિતાથી તેઓ સુલભ જગતોમાં પણ
સત્ય હોય છે. નિયમ માત્ર એ જ ઉપસર્ગ પર નિષ્કર્ષ
લખે છે એમ સ્પષ્ટ કર્યું છે.}
$\land$ અને $\lor$ના નિયમો \olref{tab:prop-rules}માં
એકત્ર કર્યા છે.

\begin{table}
  \[\def\arraystretch{3}\begin{array}{|c|c|}
    \hline
    \AxiomC{\sFmla{\True}{!A \land !B}[\sigma]}
    \RightLabel{\TRule{\True}{\land}}
    \UnaryInfC{\sFmla{\True}{!A}[\sigma]}
    \noLine
    \UnaryInfC{\sFmla{\True}{!B}[\sigma]}
    \DisplayProof
    &
    \AxiomC{\sFmla{\False}{!A \land !B}[\sigma]}
    \RightLabel{\TRule{\False}{\land}}
    \UnaryInfC{$\sFmla{\False}{!A}[\sigma] \quad \mid \quad
      \sFmla{\False}{!B}[\sigma]$}
    \DisplayProof
    \\[2ex]
    \hline
    \AxiomC{\sFmla{\True}{!A \lor !B}[\sigma]}
    \RightLabel{\TRule{\True}{\lor}}
    \UnaryInfC{$\sFmla{\True}{!A}[\sigma] \quad \mid \quad
      \sFmla{\True}{!B}[\sigma]$}
    \DisplayProof
    &
    \AxiomC{\sFmla{\False}{!A \lor !B}[\sigma]}
    \RightLabel{\TRule{\False}{\lor}}
    \UnaryInfC{\sFmla{\False}{!A}[\sigma]}
    \noLine
    \UnaryInfC{\sFmla{\False}{!B}[\sigma]}
    \DisplayProof
    \\[2ex]
    \hline
  \end{array}\]
  \caption{$\land$ અને $\lor$ માટે ઉપસર્ગવાળા ટેબ્લો નિયમો}
  \ollabel{tab:prop-rules}
\end{table}

બંધતાની શરત સામાન્ય ટેબ્લો જેવી છે, પરંતુ
સૂત્રો સાથે ઉપસર્ગો પણ સુસંગત હોવા જોઈએ. વધુમાં,
અહીં થોડી વ્યાપક શરત લઈ શકીએ: સ્ફુરણાવાદી
નિદર્શોમાં સૂત્રો એક વાર સત્ય બને, પછી
સત્ય રહે છે. તેથી $!A$ એ~$\sigma$માં સત્ય હોય
અને સુલભ ઉપસર્ગ~$\sigma.{*}$ પર ખોટું હોય
એવું શક્ય નથી. આથી કોઈ શાખામાં નીચેના બંને હોય,
તો તે બંધ છે:
\[
\sFmla{\True}{!A}[\sigma] \quad\text{અને}\quad \sFmla{\False}{!A}[\sigma.{*}]
\]
કોઈ ઉપસર્ગ $\sigma$ અને સૂત્ર~$!A$ માટે.
નોંધો કે ચિહ્નો ઊલટાં હોય, એટલે શાખામાં
\[
\sFmla{\False}{!A}[\sigma] \quad\text{અને}\quad \sFmla{\True}{!A}[\sigma.{*}]
\]
હોય, તો $*$ ખાલી શ્રેણી હોય ત્યારે જ શાખા બંધ છે.

આ ઉપરાંત, શાખામાં~$\sFmla{\True}{\bot}[\sigma]$
હોય, તો તે બંધ છે.

ધારણાઓ ગોઠવવાના નિયમો પણ સામાન્ય ટેબ્લો
જેવા જ છે, પરંતુ ધારણાઓ માટે આપણે હંમેશા
ઉપસર્ગ~$1$ લઈએ છીએ. (કયો ઉપસર્ગ લઈએ તે
મહત્ત્વનું નથી, જો બધી ધારણાઓ માટે તે એક જ હોય.)
દાખલા તરીકે, આપણે કહીએ છીએ કે
\[
!B_1, \dots, !B_n \Proves !A
\]
ત્યારે અને માત્ર ત્યારે જ જ્યારે નીચેની ધારણાઓ
માટે બંધ ટેબ્લો હોય:
\[
\sFmla{\True}{!B_1}[1], \dots, \sFmla{\True}{!B_n}[1],
\sFmla{\False}{!A}[1].
\]

શરતી~$\lif$ માટેના નિયમો પ્રશિષ્ટ અને મોડલ
કિસ્સાઓથી જુદા છે. $\TRule{\True}{\lif}$ નિયમ
$\sFmla{\True}{!A \lif !B}[\sigma]$ ધરાવતી શાખાને
જુદી બે શાખાઓ પર અનુક્રમે
$\sFmla{\False}{!A}[\sigma.{*}]$ અને
$\sFmla{\True}{!B}[\sigma.{*}]$ ઉમેરીને વિસ્તારે છે.
\sourcecorrection{OLMD-099}{મૂળ ગદ્યમાં આ નિયમનાં
બંને નિષ્કર્ષોના સત્યચિહ્ન ઊલટાં છે: ત્યાં
$\sFmla{\True}{!A}[\sigma.{*}]$ અને
$\sFmla{\False}{!B}[\sigma.{*}]$ છે, જ્યારે
નીચેની નિયામક કોષ્ટકમાં સાચાં
$\sFmla{\False}{!A}[\sigma.{*}]$ અને
$\sFmla{\True}{!B}[\sigma.{*}]$ છે. મૂળના
બંને \texttt{\string\TRule} ગદ્ય-લેબલોમાં દલીલોનો ક્રમ પણ
કોષ્ટકથી ઊલટો છે; અહીં કોષ્ટકનો ક્રમ લીધો છે.}
આ નિયમ ફક્ત એવા ઉપસર્ગ~$\sigma.{*}$ માટે
લાગુ કરી શકાય જે સંબંધિત શાખા પર \emph{પહેલેથી}
આવે છે. આવા ઉપસર્ગને તે શાખા પર ``વપરાયેલો''
કહીએ. ($\sigma.{*}$માં $\sigma$ પોતે પણ આવે છે,
તેથી જુદી શાખાઓ પર અનુક્રમે
$\sFmla{\False}{!A}[\sigma]$ અને
$\sFmla{\True}{!B}[\sigma]$ ઉમેરીને આ નિયમ
હંમેશા લાગુ કરી શકાય છે.)

$\TRule{\False}{\lif}$ નિયમ
$\sFmla{\False}{!A \lif !B}[\sigma]$ ધરાવતી શાખાને
એ જ શાખા પર $\sFmla{\True}{!A}[\sigma.n]$
અને $\sFmla{\False}{!B}[\sigma.n]$ બંને ઉમેરીને
વિસ્તારે છે; અહીં $\sigma.n$ એ શાખા પર નવો
ઉપસર્ગ છે.

$\lnot$ માટેના નિયમો સમાન રીતે વ્યાખ્યાયિત થાય છે
($\lnot !A$ને $!A \lif \lfalse$ તરીકે વ્યાખ્યાયિત
કર્યા પ્રમાણે).

નિયમો \olref{tab:rules-lif-lnot}માં આપ્યા છે.

\begin{table}
  \[\def\arraystretch{3}\begin{array}{|c|c|}
    \hline
    \AxiomC{\sFmla{\True}{\lnot !A}[\sigma]}
    \RightLabel{\TRule{\True}{\lnot}}
    \UnaryInfC{$\sFmla{\False}{!A}[\sigma.{*}]$}
    \DisplayProof
    &
    \AxiomC{\sFmla{\False}{\lnot !A}[\sigma]}
    \RightLabel{\TRule{\False}{\lnot}}
    \UnaryInfC{\sFmla{\True}{!A}[\sigma.n]}
    \DisplayProof
    \\[1ex]
    \text{$\sigma.{*}$ વપરાયેલો છે} & \text{$\sigma.n$ નવો છે}\\
    \hline
    \AxiomC{\sFmla{\True}{!A \lif !B}[\sigma]}
    \RightLabel{\TRule{\True}{\lif}}
    \UnaryInfC{$\sFmla{\False}{!A}[\sigma.{*}] \quad \mid \quad
      \sFmla{\True}{!B}[\sigma.{*}]$}
    \DisplayProof
    &
    \AxiomC{\sFmla{\False}{!A \lif !B}[\sigma]}
    \RightLabel{\TRule{\False}{\lif}}
    \UnaryInfC{\sFmla{\True}{!A}[\sigma.n]}
    \noLine
    \UnaryInfC{\sFmla{\False}{!B}[\sigma.n]}
    \DisplayProof
    \\[1ex]
    \text{$\sigma.{*}$ વપરાયેલો છે} & \text{$\sigma.n$ નવો છે}\\
    \hline
  \end{array}\]
  \caption{$\lnot$ અને $\lif$ માટે ઉપસર્ગવાળા ટેબ્લો નિયમો}
  \ollabel{tab:rules-lif-lnot}
\end{table}
% END OLP-0513
\endgroup

\begingroup
\makeatletter\def\ol@sectioncs{section}\makeatother

% BEGIN OLP-0514
\olfileid{int}{tab}{prf}

\olsection{સ્ફુરણાવાદી તર્ક માટેનું ટેબ્લો}
\phantomsection\label{guunit:OLP-0514}


\begin{ex}
  $(!A \land !B) \lif !C \Proves
  !A \lif (!B \lif !C)$ બતાવતું બંધ ટેબ્લો આપીએ છીએ.
  \begin{oltableau}
    [\pFmla{\True}{(\formula{A} \land \formula{B}) \lif \formula{C}}{1},
      just =\TAss
      [\pFmla{\False}{\formula{A} \lif (\formula{B} \lif \formula{C})}{1},
        just =\TAss
        [\pFmla{\True}{\formula{A}}{1.1},
          just = {\TRule{\False}{\lif}[2]}
          [\pFmla{\False}{\formula{B} \lif \formula{C}}{1.1},
            just = {\TRule{\False}{\lif}[2]}
            [\pFmla{\True}{\formula{B}}{1.1.1},
              just = {\TRule{\False}{\lif}[4]}
              [\pFmla{\False}{\formula{C}}{1.1.1},
                just = {\TRule{\False}{\lif}[4]}
                [\pFmla{\False}{\formula{A} \land \formula{B}}{1.1.1},
                  just = {\TRule{\True}{\lif}[1]}
                  [\pFmla{\False}{\formula{A}}{1.1.1},
                    just= {\TRule{\False}{\land}[7]}, close]
                  [\pFmla{\False}{\formula{B}}{1.1.1},
                    just= {\TRule{\False}{\land}[7]}, close]]
                [\pFmla{\True}{\formula{C}}{1.1.1},
                  just = {\TRule{\True}{\lif}[1]}, close]]
            ]
          ]
        ]
      ]
    ]
  \end{oltableau}
  \sourcecorrection{OLMD-100}{મૂળના બંધ ટેબ્લોમાં
  $\False$-સંયોજનના બે તારણો માટે આધાર-પંક્તિ [4]
  બતાવેલી છે, પરંતુ [4] એ $\False$-શરતી છે.
  બંને તારણો [7]માં આવેલા $\False$-સંયોજનમાંથી
  મળે છે; બે આધાર-સંદર્ભ [7] કર્યા છે.}
\end{ex}

\begin{prob}
  નીચેના વિધાનો બતાવવા બંધ સ્ફુરણાવાદી ટેબ્લો
  શોધો:
  \begin{enumerate}
    \item $\Proves !A \lif (!B \lif !A)$
    \item $\Proves \lnot(!A \land \lnot !A)$
    \item $!A \lif (!B \lif !C) \Proves (!A \land !B) \lif !C$
    \item $\lnot !A \lor \lnot !B \Proves \lnot(!A \land !B)$
  \end{enumerate}
\end{prob}
% END OLP-0514
\endgroup

\begingroup
\makeatletter\def\ol@sectioncs{section}\makeatother

% BEGIN OLP-0515
\olfileid{int}{tab}{sou}

\olsection{સ્ફુરણાવાદી ટેબ્લોની યથાર્થતા}
\phantomsection\label{guunit:OLP-0515}


\begin{explain}
  સ્ફુરણાવાદી ટેબ્લો યથાર્થ છે તે બતાવવા,
  જો
  \[
  \sFmla{\True}{!B_1}[1], \dots, \sFmla{\True}{!B_n}[1], \sFmla{\False}{!A}[1]
  \]
  માટે બંધ ટેબ્લો હોય, તો
  $!B_1, \dots, !B_n \Entails !A$ છે તે બતાવવું પડે.
  પ્રતિપક્ષી વિધાન સાબિત કરવું સરળ છે: જો કોઈ
  $\mModel{M}$ અને જગત~$w$ માટે બધા $i=1$,
  \dots,~$n$ માટે $\mSat{M}{!B_i}[w]$ હોય,
  પરંતુ $\mSat/{M}{!A}[w]$ હોય, તો કોઈ
  ટેબ્લો બંધ થઈ શકે નહીં.
  \sourcecorrection{OLMD-101}{મૂળમાં અહીં
  પ્રતિનિદર્શમાં $!A$ પણ સત્ય છે એમ લખાયું છે;
  $!A$ ખોટું હોવું જોઈએ, તેથી સત્ય-સંજ્ઞામાં
  નિષેધચિહ્ન ઉમેર્યું છે.}
  આવું પ્રતિનિદર્શ બતાવે છે કે ટેબ્લોની
  શરૂઆતની ધારણાઓ સંતોષ્ય છે. સાબિતીની રીત એ
  બતાવવાની છે કે ટેબ્લોની શાખા પરનાં બધાં
  ઉપસર્ગવાળા સૂત્રો સંતોષ્ય હોય, તો કોઈ પણ
  નિયમ લાગુ કરતાં ઓછામાં ઓછી એક વિસ્તરેલી શાખા
  સંતોષ્ય રહે છે. બંધ શાખાઓ અસંતોષ્ય હોવાથી,
  સંતોષ્ય ઉપસર્ગવાળા સૂત્રોના ગણનું કોઈ પણ
  ટેબ્લો ઓછામાં ઓછી એક ખુલ્લી શાખા ધરાવે છે.

  સ્ફુરણાવાદી કિસ્સામાં આ રીત લાગુ કરવા
  ``સંતોષ્ય''ની વ્યાખ્યાને સંબંધાત્મક નિદર્શો અને
  ઉપસર્ગો સુધી વિસ્તૃત કરવી પડે છે. પછી સાબિતી
  સીધી છે.
  \sourcecorrection{OLMD-102}{મૂળમાં અહીં
  ``મોડલ કિસ્સો'' લખાયો છે; આ પ્રકરણ
  સ્ફુરણાવાદી ટેબ્લોનું છે, તેથી વિષયનું નામ
  સુસંગત કર્યું છે.}
\end{explain}

\begin{defn}
  ધારો કે $P$ ઉપસર્ગોનો ગણ છે, એટલે
  $P \subseteq (\PosInt)^* \setminus \{\emptyseq\}$,
  અને $\mModel{M}$ નિદર્શ છે.
  વિધેય~$f\colon P \to W$ એ~$\mModel{M}$માં
  $P$નું \emph{અર્થઘટન} ત્યારે કહેવાય જ્યારે
  $\sigma$ અને $\sigma.n$ બંને~$P$માં હોય ત્યારે
  $Rf(\sigma)f(\sigma.n)$ હોય.

  ઉપસર્ગો~$P$ના અર્થઘટનની દૃષ્ટિએ નીચેની
  વ્યાખ્યાઓ આપી શકીએ:
  \begin{enumerate}
  \item $\mModel{M}$ એ $\sFmla{\True}{!A}[\sigma]$ને
    ત્યારે અને માત્ર ત્યારે જ સંતોષે જ્યારે
    $\mSat{M}{!A}[f(\sigma)]$.
  \item $\mModel{M}$ એ $\sFmla{\False}{!A}[\sigma]$ને
    ત્યારે અને માત્ર ત્યારે જ સંતોષે જ્યારે
    $\mSat/{M}{!A}[f(\sigma)]$.
  \end{enumerate}
\end{defn}

$R$ સ્વવાચક અને પરંપરિત હોવાથી, અને $\sigma.{*}$
એ $\sigma$, $\sigma.n_1$, $\sigma.n_1.n_2$,
\dots, દર્શાવતું હોવાથી, $Rf(\sigma)f(\sigma.{*})$
પણ છે.

\begin{defn}
  ધારો કે $\Gamma$ ઉપસર્ગવાળા સૂત્રોનો ગણ છે,
  અને $P(\Gamma)$ તેમાં આવતા ઉપસર્ગોનો ગણ છે.
  જો $f$ એ $\mModel{M}$માં $P(\Gamma)$નું અર્થઘટન
  હોય, તો $\mModel{M}$ એ~$f$ની દૃષ્ટિએ $\Gamma$ને
  સંતોષે છે, સંજ્ઞામાં $\mSat{M}{\Gamma}[f]$,
  એમ ત્યારે કહીએ જ્યારે $\mModel{M}$ એ~$f$ની
  દૃષ્ટિએ $\Gamma$નું દરેક ઉપસર્ગવાળું સૂત્ર
  સંતોષે. $\Gamma$ \emph{સંતોષ્ય} ત્યારે અને માત્ર
  ત્યારે જ જ્યારે એવું નિદર્શ~$\mModel{M}$ અને
  $P(\Gamma)$નું અર્થઘટન~$f$ હોય કે
  $\mSat{M}{\Gamma}[f]$.
\end{defn}

\begin{prop}
  જો $\Gamma$માં કોઈ સૂત્ર~$!A$ અને
  ઉપસર્ગ~$\sigma$ માટે
  $\sFmla{\True}{!A}[\sigma]$ તથા
  $\sFmla{\False}{!A}[\sigma.{*}]$ બંને હોય,
  અથવા તેમાં $\sFmla{\True}{\lfalse}[\sigma]$
  હોય, તો $\Gamma$ અસંતોષ્ય છે.
\end{prop}

\begin{proof}
  હંમેશા $\mSat/{M}{\lfalse}[f(\sigma)]$ હોવાથી,
  $\sFmla{\True}{\lfalse}[\sigma]$ ધરાવતો
  $\Gamma$ અસંતોષ્ય છે.

  વળી, બંને વિરોધી ચિહ્નિત સૂત્રોને સંતોષતું
  નિદર્શ~$\mModel{M}$ અને $P(\Gamma)$નું અર્થઘટન~$f$
  હોઈ શકે નહીં. જો $\mSat{M}{!A}[f(\sigma)]$
  હોય, તો $Rf(\sigma)f(\sigma.{*})$ અને
  \olref[sem][rel]{prop:true-monotonic} મુજબ
  $\mSat{M}{!A}[f(\sigma.{*})]$ પણ હોય.
  તેથી $\mSat{M}{!A}[f(\sigma)]$ અને
  $\mSat/{M}{!A}[f(\sigma.{*})]$ બંને સાથે
  હોઈ શકે નહીં.
  \sourcecorrection{OLMD-103}{મૂળના આ વાક્યમાં
  વાક્યરચના અધૂરી છે, $R$ની બીજી દલીલમાં
  $f$ ખૂટે છે અને એકદિશવર્ધિતાથી મળતું સત્ય
  ફરી મૂળ જગતમાં મૂકાયું છે. અહીં તે સુલભ
  જગત~$f(\sigma.{*})$માં મૂક્યું છે; અસત્યના
  ચિહ્નિત સૂત્રને પણ~$\sigma$ ઉપસર્ગ આપ્યો છે.}
\end{proof}

\begin{thm}[યથાર્થતા]
  \ollabel{thm:tableau-soundness}
  જો $\Gamma$ માટે બંધ ટેબ્લો હોય,
  તો $\Gamma$ અસંતોષ્ય છે.
\end{thm}

\begin{proof}
ટેબ્લોની શાખાને ત્યારે સંતોષ્ય કહીએ જ્યારે
તેના પરના ચિહ્નિત સૂત્રોનો ગણ સંતોષ્ય હોય;
અને ટેબ્લોને ત્યારે સંતોષ્ય કહીએ જ્યારે
તેની ઓછામાં ઓછી એક શાખા સંતોષ્ય હોય.

આ બતાવીએ: સંતોષ્ય ટેબ્લોને અનુમાનનો કોઈ પણ
નિયમ વાપરીને વિસ્તારો, તો મળતું ટેબ્લો પણ
સંતોષ્ય રહે છે. આનાથી પ્રમેય મળશે: કોઈ પણ બંધ
ટેબ્લો એ~$\Gamma$ની ધારણાઓથી જ શરૂ થતા
ટેબ્લો પર અનુમાનના નિયમો લગાવવાથી મળે છે.
જો $\Gamma$ સંતોષ્ય હોય, તો તેના માટેનું દરેક
ટેબ્લો સંતોષ્ય રહે. પરંતુ બંધ ટેબ્લો
સંતોષ્ય નથી, કારણ કે તેની બધી શાખાઓ બંધ છે
અને બંધ શાખાઓ અસંતોષ્ય છે.

ધારો કે આપણી પાસે સંતોષ્ય ટેબ્લો છે,
એટલે ઓછામાં ઓછી એક સંતોષ્ય શાખાવાળું
ટેબ્લો. નિયમ લાગુ કરતાં કોઈ શાખા પર
ચિહ્નિત સૂત્રો ઉમેરાય છે અથવા શાખા
બે ભાગમાં ફાટે છે. જો ટેબ્લોમાં નિયમથી ન વિસ્તરતી
બીજી કોઈ સંતોષ્ય શાખા હોય, તો તે વિસ્તરેલા
ટેબ્લોમાં પણ સંતોષ્ય રહે છે; તેથી આખું
વિસ્તરેલું ટેબ્લો સંતોષ્ય છે. હવે ફક્ત એવો
કિસ્સો જોવો છે જેમાં નિયમ સંતોષ્ય શાખા પર
લાગુ થાય છે.

તે શાખા પરનાં ચિહ્નિત સૂત્રોનો ગણ
$\Gamma$ હોય અને નિયમ જેના પર લાગુ થાય છે તે
ચિહ્નિત સૂત્ર $\sFmla{S}{!A}[\sigma] \in \Gamma$
હોય. નિયમથી શાખા ન ફાટે, તો વિસ્તરેલી શાખા,
એટલે નિયમના નિષ્કર્ષો સાથેનો $\Gamma$,
હજી સંતોષ્ય છે તે બતાવવું પડે. શાખા ફાટે,
તો બેમાંથી ઓછામાં ઓછી એક નવી શાખા સંતોષ્ય
છે તે બતાવવું પડે.
\begin{enumerate}
\item $\sFmla{\True}{\lnot !B}[\sigma] \in \Gamma$
  પર $\TRule{\True}{\lnot}$ લગાવવાથી શાખા
  વિસ્તરે છે. વિસ્તરેલી શાખામાં
  ચિહ્નિત સૂત્રોનો ગણ $\Gamma \cup
  \{\sFmla{\False}{!B}[\sigma.{*}]\}$ હોય છે.
  ધારો કે $\mSat{M}{\Gamma}[f]$. ખાસ કરીને
  $\mSat{M}{\lnot !B}[f(\sigma)]$. તેથી
  $Rf(\sigma)w$ ધરાવતા કોઈ પણ $w$ માટે
  $\mSat/{M}{!B}[w]$; તેમાં $f(\sigma.{*})$
  પણ આવે છે. આથી $\mModel{M}$ એ~$f$ની
  દૃષ્ટિએ $\sFmla{\False}{!B}[\sigma.{*}]$
  સંતોષે છે.
\item $\sFmla{\False}{\lnot !B}[\sigma] \in \Gamma$
  પર $\TRule{\False}{\lnot}$ લગાવવાથી શાખા
  વિસ્તરે છે: સ્વાધ્યાય.
\item $\sFmla{\True}{!B \land !C}[\sigma] \in \Gamma$
  પર $\TRule{\True}{\land}$ લગાવવાથી શાખા
  વિસ્તરે છે અને તે જ શાખા પર બે નવાં
  ચિહ્નિત સૂત્રો, $\sFmla{\True}{!B}[\sigma]$
  અને $\sFmla{\True}{!C}[\sigma]$, મળે છે.
  ધારો કે $\mSat{M}{\Gamma}[f]$; ખાસ કરીને
  $\mSat{M}{!B \land !C}[f(\sigma)]$.
  ત્યારે $\mSat{M}{!B}[f(\sigma)]$ અને
  $\mSat{M}{!C}[f(\sigma)]$. એટલે
  $\mModel{M}$ એ~$f$ની દૃષ્ટિએ
  $\sFmla{\True}{!B}[\sigma]$ અને
  $\sFmla{\True}{!C}[\sigma]$ બંને સંતોષે છે.
\item $\sFmla{\False}{!B \lor !C}[\sigma] \in \Gamma$
  પર $\TRule{\False}{\lor}$ લગાવવાથી શાખા
  વિસ્તરે છે: સ્વાધ્યાય.
  \sourcecorrection{OLMD-104}{મૂળની આ
  સ્વાધ્યાય-પંક્તિમાં ચિહ્નિત વિયોજન પછી
  ઉપસર્ગ~$\sigma$ ખૂટે છે; નિયમનો વિષય
  નિશ્ચિત કરવા તે ઉમેર્યો છે.}
\item $\sFmla{\False}{!B \lif !C}[\sigma] \in \Gamma$
  પર $\TRule{\False}{\lif}$ લગાવવાથી શાખા
  વિસ્તરે છે. એ જ શાખા પર બે નવાં
  ચિહ્નિત સૂત્રો મળે છે:
  $\sFmla{\True}{!B}[\sigma.n]$ અને
  $\sFmla{\False}{!C}[\sigma.n]$; અહીં
  $\sigma.n$ શાખા પર નવો ઉપસર્ગ છે, એટલે
  $\sigma.n \notin P(\Gamma)$.

  $\Gamma$ સંતોષ્ય હોવાથી એવું
  $\mModel{M}$ અને $P(\Gamma)$નું અર્થઘટન~$f$
  છે કે $\mSat{M}{\Gamma}[f]$; ખાસ કરીને
  $\mSat/{M}{!B \lif !C}[f(\sigma)]$.
  $\Gamma \cup \{\sFmla{\True}{!B}[\sigma.n],
  \sFmla{\False}{!C}[\sigma.n]\}$ સંતોષ્ય છે
  તે બતાવવાનું છે. તે માટે $P(\Gamma)
  \cup \{\sigma.n\}$નું અર્થઘટન નીચે મુજબ
  વ્યાખ્યાયિત કરીએ છીએ:
  \sourcecorrection{OLMD-105}{મૂળ આ લક્ષ્યમાં
  બે નવા ચિહ્નિત સૂત્રોને બદલે
  $\sFmla{\False}{!B \lif !C}[\sigma.n]$
  લખે છે; એ શાખામાં લાગુ પડેલા નિયમનું
  નિષ્કર્ષ નથી. બંને જરૂરી સૂત્રો મૂક્યા છે.}

  $\mSat/{M}{!B \lif !C}[f(\sigma)]$
  હોવાથી એવું $w \in W$ છે કે $Rf(\sigma)w$
  અને $\mSat{M}{!B}[w]$ તથા
  $\mSat/{M}{!C}[w]$. $f'$ એ $f$ જેવું જ હોય,
  ફક્ત $f'(\sigma.n) = w$ મૂકો.
  $f'(\sigma) = f(\sigma)$ અને $Rf(\sigma)w$
  હોવાથી $Rf'(\sigma)f'(\sigma.n)$; તેથી
  $f'$ એ $P(\Gamma) \cup \{\sigma.n\}$નું
  અર્થઘટન છે. સ્પષ્ટ છે કે
  $\mSat{M}{!B}[f'(\sigma.n)]$ અને
  $\mSat/{M}{!C}[f'(\sigma.n)]$.
  $\sigma' \in P(\Gamma)$ એવા દરેક ઉપસર્ગ માટે
  $f(\sigma') = f'(\sigma')$ હોવાથી
  $\mSat{M}{\Gamma}[f']$. તેથી $\mModel{M}, f'$
  એ $\Gamma \cup \{\sFmla{\True}{!B}[\sigma.n],
  \sFmla{\False}{!C}[\sigma.n]\}$ સંતોષે છે.
\end{enumerate}
હવે બે નવી શાખાઓ આપતા નિયમો જોઈએ.
\begin{enumerate}
\item $\sFmla{\False}{!B \land !C}[\sigma] \in \Gamma$
  પર $\TRule{\False}{\land}$ લગાવવાથી
  શાખા બેમાં ફાટે છે: ડાબી શાખા
  $\sFmla{\False}{!B}[\sigma]$થી આગળ વધે છે
  અને જમણી $\sFmla{\False}{!C}[\sigma]$થી.
  ધારો કે $\mSat{M}{\Gamma}[f]$; ખાસ કરીને
  $\mSat/{M}{!B \land !C}[f(\sigma)]$.
  ત્યારે $\mSat/{M}{!B}[f(\sigma)]$ અથવા
  $\mSat/{M}{!C}[f(\sigma)]$. પહેલા કિસ્સામાં
  $\mModel{M}, f$ એ $\sFmla{\False}{!B}[\sigma]$
  સંતોષે છે, એટલે ડાબી શાખા સંતોષ્ય છે.
  બીજા કિસ્સામાં $\mModel{M}, f$ એ
  $\sFmla{\False}{!C}[\sigma]$ સંતોષે છે,
  એટલે જમણી શાખા સંતોષ્ય છે.
\item $\sFmla{\True}{!B \lor !C}[\sigma] \in \Gamma$
  પર $\TRule{\True}{\lor}$ લગાવવાથી શાખા
  વિસ્તરે છે: સ્વાધ્યાય.
\item $\sFmla{\True}{!B \lif !C}[\sigma] \in \Gamma$
  પર $\TRule{\True}{\lif}$ લગાવવાથી શાખા
  વિસ્તરે છે: સ્વાધ્યાય.
\end{enumerate}
\end{proof}

\begin{prob}
\olref[int][tab][sou]{thm:tableau-soundness}ની
સાબિતી પૂરી કરો.
\end{prob}

\begin{cor}
\ollabel{cor:entailment-soundness}
જો $\Gamma \Proves !A$ હોય, તો
$\Gamma \Entails !A$.
\end{cor}

\begin{proof}
  જો $\Gamma \Proves !A$ હોય, તો કોઈ
  $!B_1$, \dots, $!B_n \in \Gamma$ માટે
  $\Delta = \{\sFmla{\False}{!A}[1],
  \sFmla{\True}{!B_1}[1], \dots,
  \sFmla{\True}{!B_n}[1]\}$નું બંધ
  ટેબ્લો છે. $\Gamma \Entails !A$
  બતાવવું છે. ઊલટું ધારો: કોઈ
  $\mModel{M}$ અને $w$ માટે $i=1$,
  \dots,~$n$ માટે $\mSat{M}{!B_i}[w]$
  હોય, પણ $\mSat/{M}{!A}[w]$.
  $f(1) = w$ મૂકો; ત્યારે $f$ એ
  $P(\Delta)$નું~$\mModel{M}$માં અર્થઘટન
  છે અને $\mModel{M}$ એ~$f$ની દૃષ્ટિએ
  $\Delta$ને સંતોષે છે. પરંતુ
  \olref{thm:tableau-soundness} મુજબ બંધ
  ટેબ્લો ધરાવતો $\Delta$ અસંતોષ્ય છે;
  આ વિરોધાભાસ છે. તેથી
  $\Gamma \Entails !A$ હોવું જ જોઈએ.
  \sourcecorrection{OLMD-106}{મૂળની અંતિમ
  પંક્તિમાં ફરી $\Gamma \Proves !A$ લખાયું છે,
  જે પહેલેથી ધારેલું હતું. વિરોધાભાસથી
  અહીં જોઈતું $\Gamma \Entails !A$ મળે છે.}
\end{proof}

\begin{cor}
\ollabel{cor:weak-soundness}
જો $\Proves !A$ હોય, તો બધા નિદર્શોમાં
$!A$ સત્ય છે.
\end{cor}
% END OLP-0515
\endgroup


\OLEndChapterHook
% END OLP-0511
\endgroup


\OLEndPartHook
% END OLP-0491
